\documentclass[11pt]{article}
\usepackage[T1]{fontenc}
\usepackage[utf8]{inputenc}
\usepackage{lmodern}
\usepackage{microtype}
\usepackage[margin=1in]{geometry}
\usepackage{amsmath,amssymb,amsthm,mathtools}
\usepackage{booktabs,longtable,array,tabularx}
\usepackage{enumitem}
\usepackage{xcolor}
\usepackage{listings}
\usepackage{xurl}
\usepackage{hyperref}
\usepackage{fancyhdr}
\usepackage{lastpage}
\usepackage{graphicx}

\hypersetup{
  colorlinks=true,
  linkcolor=blue!45!black,
  citecolor=blue!45!black,
  urlcolor=blue!55!black,
  pdftitle={Certified Zariski Descent in Intuitionistic Primitive Recursive Arithmetic},
  pdfauthor={OpenAI, prepared for Vladimir Reshetnikov}
}
\urlstyle{same}
\setlist{nosep,leftmargin=2em}
\setlength{\parindent}{0pt}
\setlength{\parskip}{0.55em}
\setlength{\headheight}{14pt}
\setlength{\emergencystretch}{2.5em}
\pagestyle{fancy}
\fancyhf{}
\lhead{Certified Zariski descent in iPRA}
\rhead{\thepage/\pageref{LastPage}}
\cfoot{}

\newtheorem{theorem}{Theorem}[section]
\newtheorem{lemma}[theorem]{Lemma}
\newtheorem{proposition}[theorem]{Proposition}
\newtheorem{corollary}[theorem]{Corollary}
\newtheorem{definition}[theorem]{Definition}
\newtheorem{warning}[theorem]{Warning}
\newtheorem{remark}[theorem]{Remark}
\newtheorem{problem}[theorem]{Formalization target}

\newcommand{\N}{\mathbb N}
\newcommand{\Z}{\mathbb Z}
\newcommand{\Q}{\mathbb Q}
\newcommand{\Spec}{\operatorname{Spec}}
\newcommand{\Proj}{\operatorname{Proj}}
\newcommand{\im}{\operatorname{im}}
\newcommand{\cO}{\mathcal O}
\newcommand{\Prf}{\mathsf{Prf}}
\newcommand{\PRAAx}{\mathsf{PRAAx}}
\newcommand{\PR}{\mathsf{PR}}
\newcommand{\PTm}{\mathsf{PTm}}
\newcommand{\Fm}{\mathsf{Fm}}
\newcommand{\EqW}{\mathsf{EqW}}
\newcommand{\Valid}{\mathsf{Valid}}
\newcommand{\code}[1]{\nolinkurl{#1}}
\newcommand{\true}{\mathsf{true}}
\newcommand{\false}{\mathsf{false}}
\newcommand{\ceil}[1]{\left\lceil #1\right\rceil}
\newcommand{\abs}[1]{\left|#1\right|}
\newcommand{\set}[1]{\left\{#1\right\}}

\lstset{
  basicstyle=\ttfamily\small,
  columns=fullflexible,
  breaklines=true,
  frame=single,
  xleftmargin=0.5em,
  xrightmargin=0.5em,
  showstringspaces=false,
  keywordstyle=\bfseries,
  commentstyle=\itshape\color{black!65}
}

\title{\textbf{Certified Zariski Descent in Intuitionistic Primitive Recursive Arithmetic}\\[0.5em]
\large Finite presentations, localization, affine atlases, and projective-space cohomology}
\author{Prepared for Vladimir Reshetnikov}
\date{15 November 2026}

\begin{document}
\maketitle

\begin{abstract}
The existing \emph{ProveIt} iPRA packets develop finite contracts, demand-driven observations,
relational finite-use composition, stable outputs, exact arithmetic, and several Diophantine
constructions.  This report adds a specifically algebraic-geometric layer.  Mathematical objects
are represented by natural-number codes with proof-relevant equality certificates; quotient rings
use explicit ideal-combination witnesses; localizations use explicit annihilating exponents; and
finite affine covers are carried by unit-ideal certificates rather than by quantification over prime
ideals.  The central result is an explicit primitive-recursive Zariski gluing algorithm for modules
on a finite principal-open cover.  It works over rings with zero divisors and separates certificate
verification from certificate production.

This provides a finitary semantics for affine schemes, finite affine atlases, quasi-coherent module
data, vector bundles, and finite \v{C}ech complexes.  As a developed case study, the two-chart atlas
of $\mathbf P^1$ gives a primitive-recursive normal form for the \v{C}ech cohomology of every
$\cO(d)$.  The construction extends to the standard cover of $\mathbf P^n$: the complex decomposes
by exponent vectors, mixed-sign summands have an explicit contracting homotopy, and only the
all-nonnegative and all-negative summands survive.  Thus the familiar cohomology table is reduced
to finite combinatorics suitable for iPRA.

The archive contains complete human-readable proofs, exact Python validation, source-level Lean
interfaces, and a dependency-ordered route to actual \code{Prf PRAAx} terms.  The Lean outlines
are not claimed compiler-checked, and no general effective Siegel, Faltings, Mordell--Weil, or
singular Puiseux theorem is claimed.
\end{abstract}

\tableofcontents
\newpage

\section{Scope, thesis, and proof status}

\subsection{What the requested formalization should mean}

The phrase ``formalize algebraic geometry in iPRA'' can be understood in several inequivalent
ways.  The strongest and least useful interpretation would demand a single natural number coding
an arbitrary scheme, an internal set of all of its points, and a universal evaluator for every
algorithm mentioned in a textbook.  That interpretation is neither necessary nor compatible with
the foundational discipline of primitive recursive arithmetic.

The interpretation adopted here is presentation-relative and theorem-relative:

\begin{enumerate}[label=\arabic*.]
\item Fix a class of finite codes: polynomials, matrices, finite presentations, local fractions,
      overlap maps, and certificates.
\item Fix each primitive-recursive operation before it appears in a theorem.  There is no evaluator
      for arbitrary encoded programs.
\item Express equality and side conditions through finite witnesses whose checkers are primitive
      recursive.
\item Replace an infinite geometric object by the finite algebraic presentation needed by the
      theorem: an affine presentation, a finite atlas, a finite cover, a finite \v{C}ech complex, or
      finitely many observations of a represented object.
\item Prove the finite correctness statement with quantifier-free induction and the ordinary
      intuitionistic logical rules of iPRA.
\end{enumerate}

A completed formalization is not merely a Lean theorem about standard natural numbers.  In the
repository's notation, its endpoint is a term
\[
                 \Prf\ \PRAAx\ []\ \varphi,
\]
where $\varphi$ is the coded first-order formula expressing the advertised result.  The existing
bridge \code{pra_iff_native} may then transport that derivation to the concise native calculus.

\subsection{The central thesis}

The topological point-set definition of a scheme quantifies over prime ideals and arbitrary open
sets.  Most concrete arguments, however, consume much less information.  They use finitely
presented rings, standard opens $D(f)$, localizations, finitely many overlap maps, and finitely many
compatibility equations.  This finite algebraic skeleton is enough for a large body of algebraic
geometry.

The central thesis is therefore:

\begin{quote}
\emph{For iPRA, the right primitive notion is not an internally completed prime spectrum.  It is a
certified finite affine atlas together with proof-relevant localization and descent data.}
\end{quote}

The conventional scheme is the external mathematical interpretation of that atlas.  The internal
arithmetic proof manipulates only its finite codes and certificates.  This is analogous to working
with a finite presentation of a group rather than with a completed set containing all group
elements, but it additionally records the local-to-global structure that distinguishes geometry
from bare commutative algebra.

\subsection{Main contributions}

The report develops the following chain.

\begin{enumerate}[label=\textbf{C\arabic*.}]
\item \textbf{Certified setoids.}  Equality is a primitive-recursive relation
      $\EqW(x,y,w)$ on two object codes and a witness code.  Reflexivity, symmetry,
      transitivity, and operation congruence are implemented by witness transformers.  No equality
      search is assumed.
\item \textbf{Presented algebra.}  Equality in
      $\Z[X_0,\ldots,X_{n-1}]/(g_0,\ldots,g_{r-1})$ is witnessed by polynomials
      $h_i$ satisfying $p-q=\sum_i h_i g_i$.  Morphisms are polynomial tuples with relation
      certificates.
\item \textbf{Localization.}  Equality of $a/f^n$ and $b/f^m$ is witnessed by an exponent $k$
      and a base-ring certificate for
      $f^k(f^m a-f^n b)=0$.  The report gives explicit witness formulas for transitivity and ring
      congruence.
\item \textbf{Power covers.}  From $\sum_i c_i f_i=1$ and $N\geq1$, a bounded multinomial
      construction produces $d_i$ with $\sum_i d_i f_i^N=1$.
\item \textbf{Zariski gluing.}  Compatible local fractions in $M_{f_i}$ are homogenized, their
      overlap annihilators are synchronized, and one explicit linear combination produces a global
      section.  The same power-cover lemma proves uniqueness.
\item \textbf{Finite affine atlases.}  Charts, principal overlaps, localization isomorphisms, and
      cocycles form a finite presentation of a scheme.  Quasi-coherent modules and vector bundles
      become finite descent data.
\item \textbf{Projective geometry.}  The standard affine charts of $\Proj$ are explicit
      localizations.  For $\mathbf P^1$, Laurent-polynomial support gives complete normal forms for
      the \v{C}ech groups of $\cO(d)$.  For $\mathbf P^n$, exponent-vector summands admit an
      explicit finite contracting homotopy.
\end{enumerate}

\subsection{What is and is not proved}

The ordinary mathematical proofs of the power-cover lemma, localization calculus, principal-cover
gluing, and the projective \v{C}ech computations are complete in this report.  Their algorithms are
implemented over exact rings in the Python companion and checked on finite test suites, including a
zero-divisor example over $\Z/12\Z$.

The supplied Lean files are blueprints.  They contain host-level structures, definitions, and named
target propositions, together with the exact object-theory endpoint.  The preparation environment
had no Lean executable or local ProveIt dependency closure, so no new compilation claim is made.
The final conversion of the principal theorems to \code{Prf PRAAx} remains an engineering task,
although the required primitive recursions and quantifier-free invariants are specified.

The framework does not manufacture a general height bound, a complete rational-point list, a set of
Mordell--Weil quotient representatives, or an exact branch selector for arbitrary Cauchy
coefficients.  It makes such data composable and verifiable when supplied; it does not erase the
mathematics needed to produce them.

\section{Repository audit and the missing geometric layer}

\subsection{Pinned repository state}

The mathematical audit is pinned to commit
\code{bce5113f44d613f5625c274128ee761f35486fc5}.  During final verification, \code{main} had advanced to
\code{a2f72b0cdf57f8a3fefd34b272c3d471385e819c}, but a direct contents-API comparison showed that the three entries under
\code{Logic/PrimitiveRecursiveArithmetic} were unchanged: \code{ConciseIPRA} had Git object
\code{622927389849d1e4c3b5f813b8481b7674974acc}, \code{README.md} had object
\code{1340d58e7f9c97cc8593341f5e80452633a2b77d}, and \code{docs} had object
\code{996363773c5e9f09fa826e0cc943f55af059d8dd} at both commits.  Thus unrelated repository activity did not alter the
audited iPRA subtree.

At the pinned snapshot the root iPRA README listed sixty-six research packets.  It also stated an
important caveat: the packets were filed as delivered and were not independently rerun, compiled,
or verified by the indexing commit.  Consequently this report treats packet claims as research
inputs, not as automatically certified facts.

The checked foundation is \code{ConciseIPRA}.  The relevant Lean layers are:

\begin{center}
\begin{tabularx}{0.94\textwidth}{@{}lX@{}}
\toprule
Module/interface & Role \\
\midrule
\code{Syntax.lean} & First-order terms and formulas, substitution, and quantifier-free formulas. \\
\code{Deduction.lean} & Intuitionistic natural deduction with equality; derivations \code{Prf}. \\
\code{PR.lean} & The inductive family \code{PR}, terms \code{PTm}, axioms \code{PRAAx}, and the quantifier-free induction interface. \\
\code{NativeBiInterpretation.lean} & The bridge \code{pra_iff_native}. \\
\bottomrule
\end{tabularx}
\end{center}

This package therefore targets the usual iPRA calculus first.  Native export is the last step, not a
substitute for an embedded derivation.

\subsection{Prior frameworks that are treated as baseline}

The repository already contains a substantial general methodology.  The following table records the
parts most relevant here.

\begin{longtable}{@{}p{0.25\textwidth}p{0.31\textwidth}p{0.36\textwidth}@{}}
\toprule
Packet & Existing contribution & Boundary retained here \\
\midrule
\endfirsthead
\toprule
Packet & Existing contribution & Boundary retained here \\
\midrule
\endhead
\code{ipra_framework} & Finite contracts, bounded Booleanization, bounded search, quotient-free finite presentations, and a compiler outline. & A contract does not by itself provide geometric overlap objects, localization equality, or a sheaf gluing formula. \\
\code{ipra_embedding_framework} & Demand-driven finite observations, certificate-relative interpretation, finite windows, simple Cauchy branches, and finite algebra interfaces. & Exact global branch data and general height bounds remain separate producer problems. \\
\code{ipra_finite_observation} & Precise diagnosis of illegal universal evaluation, Cauchy-induction mistakes, computation/proof horizons, and exact-output demands. & All algorithms below are fixed PR constructions; no universal interpreter is introduced. \\
\code{ipra_coherent_methods} & Coherent finite presentations and several small genuine \code{Prf PRAAx} derivations. & The weighted and object-specific applications are not automatically embedded. \\
\code{ipra_local_global} & Variable-length recursion and numerical local-to-global rules. & Zariski localization and module descent require additional algebraic witness calculus. \\
\code{ipra_finite_use_framework} & Relational finite-use composition, backward viability, branch-invariant outputs, robust covers, and root clusters. & The present package uses that philosophy but adds exact localization, unit-ideal covers, and sheaf equalizers. \\
\code{ipra_stable_observables} & Exact or stable observables, matrix certificates, and impossibility of complete effective bound certificates in general. & A certificate verifier is not a certificate producer; this distinction remains explicit. \\
\code{diophantine_geometry_ipra} & Rational projective point codes, heights, conics, local obstructions, explicit descent, and conditional finite generation. & It treats points and heights rather than a general scheme/atlas/sheaf layer. \\
\bottomrule
\end{longtable}

The most recent finite-use packet is especially close in spirit.  It allows a producer to return a
finite relation of locally valid answers and lets a bounded consumer choose a viable transcript.
That is useful for chart selection and root clusters.  It does not, however, supply the algebraic
identity that makes a principal-open cover jointly monic or the formula that glues compatible module
fractions.  Those are the central additions here.

\subsection{Difficulties recorded by theorem-specific packets}

The theorem packets identify several logically distinct obstacles.

\begin{enumerate}[label=\arabic*.]
\item \textbf{No universal PR evaluator.}  An inductive Lean datatype of PR programs does not give
      an iPRA function evaluating arbitrary program codes.  Every evaluator in this package is fixed
      syntactically, or a finite trace is supplied and checked.
\item \textbf{No induction on infinitary predicates.}  Cauchy convergence, sheafhood over all opens,
      or ideal equality by unbounded search cannot be the induction formula.  Finite prefix folds and
      explicit witnesses replace them.
\item \textbf{Exact choice can be discontinuous.}  Singular Puiseux branches and arbitrary Cauchy
      coefficient zero tests remain outside a uniform exact interface.  Relational chart data or
      supplied separation witnesses are required.
\item \textbf{Cardinality does not locate points.}  The Faltings and Siegel packets stress that a
      bound on the number of rational or integral points is not a height cutoff and does not terminate
      enumeration.
\item \textbf{Descent requires actual representatives.}  The Mordell--Weil packet can consume a
      finite list covering a quotient but does not manufacture rational representatives from abstract
      Selmer data.
\item \textbf{Native Lean is not object iPRA.}  Several packets check host-level finite arithmetic
      but do not construct the corresponding \code{Prf PRAAx} value.
\end{enumerate}

The geometric framework below resolves none of the genuinely missing height or representative
producer theorems.  It does resolve a different structural difficulty: how to represent quotient
rings, localizations, finite covers, gluing, and projective cohomology without silently using
classical quotient types, prime spectra, arbitrary choice, or induction beyond iPRA.

\subsection{Why the Wikipedia checklist is useful but insufficient}

The user-specified overview divides algebraic geometry into polynomial zero sets, affine and
projective varieties, computational methods, and the modern scheme/sheaf viewpoint.  Each item
requires a different iPRA representation:

\begin{itemize}
\item polynomial zero sets use finite polynomial evaluation and bounded point searches;
\item affine geometry uses finitely presented coordinate rings;
\item projective geometry uses homogeneous presentations and finite standard charts;
\item computational geometry uses certificate-checked reductions rather than an assumed terminating
      universal algorithm;
\item schemes and sheaves use finite affine atlases and descent equalizers rather than an internally
      completed topology of prime ideals.
\end{itemize}

The rest of the report supplies one common architecture for these representations.

\section{A general framework: certified presentations and finite descent}

\subsection{The five-layer architecture}

A reusable embedding of mathematics into iPRA should make the transition from ordinary mathematical
language to arithmetic explicit.  The following five layers are sufficient for the examples in this
report and subsume the earlier finite-contract methodology.

\begin{description}[style=nextline]
\item[Layer 1: numerical representation.]
Every datum is a natural-number code.  Lists, trees, signed integers, polynomials, matrices, and
finite diagrams have primitive-recursive constructors and bounded accessors.  Malformed inputs are
handled by total default values; a separate predicate $\Valid(x)$ states the intended domain.

\item[Layer 2: proof-relevant equality.]
An object code need not have a canonical normal form.  Equality is witnessed by a code $w$ satisfying
$\EqW(x,y,w)$.  The checker is primitive recursive, and fixed primitive-recursive transformers
produce witnesses for reflexivity, symmetry, transitivity, and congruence.

\item[Layer 3: static algebra and finite diagrams.]
Operations are fixed PR functions on codes.  A finite diagram consists of finitely many objects,
arrows, and compatibility witnesses.  There is no quantification over an open-ended universe of
programs or objects.

\item[Layer 4: observations, covers, and descent.]
A theorem asks only for a finite observation: a bounded coefficient identity, a local section on a
finite cover, a finite \v{C}ech class, or a finite error bound.  Local answers are composed using
explicit compatibility and gluing certificates.

\item[Layer 5: object-theory reflection.]
The finite checker and constructor are represented by $\PR$ terms.  Their loop invariants are
quantifier-free formulas.  The result is assembled as a term of type $\Prf\ \PRAAx\ \Gamma\ \varphi$.
\end{description}

The architecture is deliberately asymmetric.  A witness may certify a correct output even when no
uniform procedure can discover that witness from the raw input.  This is not a defect: it isolates
the actual mathematical producer theorem.

\subsection{Certified setoids}

\begin{definition}[Certified setoid]
A certified setoid consists of fixed primitive-recursive functions
\[
\Valid(x),\qquad \EqW(x,y,w),\qquad
\mathsf{refl}(x),\quad \mathsf{sym}(x,y,w),\quad
\mathsf{trans}(x,y,z,u,v),
\]
together with iPRA derivations of
\begin{align*}
  \Valid(x)&\to \EqW(x,x,\mathsf{refl}(x)),\\
  \EqW(x,y,w)&\to \EqW(y,x,\mathsf{sym}(x,y,w)),\\
  \EqW(x,y,u)\wedge\EqW(y,z,v)
    &\to \EqW(x,z,\mathsf{trans}(x,y,z,u,v)).
\end{align*}
The existential relation $x\simeq y :\Leftrightarrow \exists w\,\EqW(x,y,w)$ is the represented
equality.
\end{definition}

No decision procedure for $x\simeq y$ is required.  A theorem using an equality carries the witness
it needs.  This representation is stronger than merely declaring an equivalence relation: it makes
proof composition an explicit finite operation.

\begin{definition}[Certified map]
A certified map $F:A\to B$ is a fixed PR function on codes, a domain-preservation proof, and a PR
witness transformer $F_*$ satisfying
\[
   \EqW_A(x,y,w)\to \EqW_B(Fx,Fy,F_*(x,y,w)).
\]
\end{definition}

Certified maps compose by ordinary function composition and composition of witness transformers.
Identity maps use the supplied reflexivity witness.  Thus represented objects and certified maps
form a category at the metalevel, while every individual composition is a finite arithmetic
construction.

\begin{proposition}[Finite extensional compilation]
Fix certified setoids, certified maps, and a finite expression tree built from their operations.  A
witness that corresponding leaves are equal can be transformed primitive-recursively into a witness
that the two expression values are equal.  The correctness proof uses only structural recursion on
the fixed expression and quantifier-free induction for the bounded list folds occurring in the
code.
\end{proposition}

\begin{proof}
At each operation node apply its congruence transformer.  At a finite-list node fold the binary
transformer while maintaining the quantifier-free invariant that the current accumulator pair has a
valid equality witness.  Since the expression tree is fixed before the theorem is formed, the
metalevel structural recursion produces a finite derivation; it is not an object-level truth
predicate.
\end{proof}

\subsection{Algebraic structures without quotient types}

A certified ring is a certified setoid equipped with PR operations $0,1,+,-,\cdot$ and witness
transformers for congruence.  Its ring laws are finite families of equality witnesses.  A certified
module similarly supplies addition, negation, scalar multiplication, and their congruence data.

This formulation accommodates several common representations:

\begin{itemize}
\item rational numbers as numerator/positive-denominator pairs with cross-multiplication witnesses;
\item projective points as primitive coordinate tuples with scaling witnesses;
\item quotient rings as polynomial codes with ideal-combination witnesses;
\item localization fractions with annihilator witnesses;
\item Cauchy names observed only to a finite precision, with quantitative closeness witnesses;
\item chain complexes whose equality is coefficientwise certificate equality.
\end{itemize}

The key gain is uniformity.  Once an operation has a congruence transformer, all later constructions
can use it without deciding a quotient equality or selecting a canonical representative.

\subsection{Finite diagrams and descent objects}

A finite diagram has a natural-number bound $r$ on its objects and finitely coded arrows.  Source,
target, identity, and composition are fixed PR functions.  Equality of arrows is proof-relevant.
A diagram-valued construction supplies object and arrow codes together with finite commutativity
witnesses.

For geometry, a cover is represented by a finite index set $i<r$, local objects $A_i$, overlap
objects $A_{ij\alpha}$, and restriction maps.  A descent datum is a tuple $s_i$ with equality
witnesses between its restrictions on each overlap, plus cocycle witnesses where needed.  A gluing
contract consists of
\begin{enumerate}[label=(\alph*)]
\item a PR constructor producing a candidate global code;
\item PR constructors for the restriction-equality witnesses;
\item a PR uniqueness transformer taking local equality witnesses to a global one.
\end{enumerate}

The principal-open case will supply all three explicitly.  More general finite sites can take them
as verified primitives or reduce them to principal refinements.

\subsection{Exact, relational, and approximate outputs}

The framework has three output modes.

\begin{description}[style=nextline]
\item[Exact certificate mode.]
A quotient identity, localization equality, or matrix equation is returned with an exact finite
witness.  This is the main mode of the algebraic geometry developed below.

\item[Relational finite-use mode.]
A producer supplies a finite relation of possible charts, branches, or refinements.  A bounded
consumer uses the repository's finite-use viability construction.  The final output must be
branch-invariant or quantitatively stable.

\item[Approximation mode.]
A represented infinite object is queried only to finitely many precisions.  A modulus, apartness
witness, or error DAG determines the demanded input precision.  Exact algebraic conclusions may not
be inferred without exact separation data.
\end{description}

These modes can be mixed.  For example, a singular variety may be covered by a finite relation of
candidate charts, while each chosen chart carries exact polynomial and localization certificates.
The finite-use framework handles chart choice; the certified atlas handles algebraic gluing.

\subsection{The no-universal-evaluator discipline}

Nothing above licenses a PR function that evaluates arbitrary encoded PR programs.  There are two
sound implementation patterns:

\begin{enumerate}[label=\arabic*.]
\item the theorem fixes a particular function code built from the constructors of $\PR$;
\item the input contains a finite execution trace and a fixed checker verifies each transition.
\end{enumerate}

A degree parameter, number of variables, cover size, or matrix dimension may be an ordinary numerical
input to one fixed evaluator.  What is forbidden is replacing that evaluator by a universal
interpreter whose termination or primitive recursiveness has not been established.

\section{Finitely presented algebra with proof-relevant equality}

\subsection{Polynomial codes}

Fix a number $n$ of variables.  A monomial is a vector
$\alpha=(\alpha_0,\ldots,\alpha_{n-1})\in\N^n$, encoded by a finite tuple.  A polynomial over
$\Z$ is a finite list of coefficient--monomial pairs.  The code need not be normalized: duplicate
monomials and zero coefficients are allowed.  Evaluation, addition, negation, multiplication,
substitution, coefficient collection on a supplied finite degree box, and equality of collected
coefficient vectors are primitive recursive.

There are two legitimate equality interfaces.

\begin{enumerate}[label=\arabic*.]
\item A canonical normalization algorithm may be used when coefficient equality is decidable and the
      required sorting/collection procedure has been represented.
\item More generally, a coefficientwise identity is supplied as a finite certificate.  This avoids
      requiring a global normal form and extends directly to coefficients in a certified ring.
\end{enumerate}

The second interface is foundational below.  Gröbner bases can accelerate witness production, but
are not required for the definition of quotient equality.

\subsection{Quotient rings by ideal-combination witnesses}

Let
\[
 A=\Z[X_0,\ldots,X_{n-1}]/(g_0,\ldots,g_{r-1}).
\]
An element of $A$ is represented by a polynomial code.  A witness that $p$ and $q$ represent the
same element is a vector of polynomial codes $h=(h_0,\ldots,h_{r-1})$ satisfying
\begin{equation}\label{eq:ideal-witness}
        p-q=\sum_{i<r} h_i g_i
\end{equation}
coefficientwise.

\begin{theorem}[Exact certified presentation of a quotient ring]\label{thm:quotient-setoid}
The checker for \eqref{eq:ideal-witness} and the following witness transformers make polynomial
codes into a certified commutative ring presenting $A$:
\begin{align*}
\mathsf{refl}(p)&=(0,\ldots,0),\\
\mathsf{sym}(h)&=(-h_0,\ldots,-h_{r-1}),\\
\mathsf{trans}(h,k)&=(h_0+k_0,\ldots,h_{r-1}+k_{r-1}).
\end{align*}
Addition, negation, and multiplication admit primitive-recursive congruence transformers.
\end{theorem}

\begin{proof}
Reflexivity, symmetry, and transitivity follow by the corresponding polynomial identities.  Suppose
$p-p'=\sum h_i g_i$ and $q-q'=\sum k_i g_i$.  Then
\[
 (p+q)-(p'+q')=\sum_i(h_i+k_i)g_i,
\]
and
\[
 pq-p'q'=(p-p')q+p'(q-q')
         =\sum_i(h_iq+p'k_i)g_i.
\]
These formulas are the required witness transformers.  Negation sends $h_i$ to $-h_i$.  Ring laws
are literal polynomial identities and therefore have fixed zero or elementary coefficient
certificates.  Every operation is a bounded fold over finite polynomial lists, so its correctness is
proved by quantifier-free induction on a list index.
\end{proof}

The relation is complete by definition of the generated ideal: $p$ and $q$ have equal quotient
classes exactly when some finite list $h_i$ satisfies \eqref{eq:ideal-witness}.  Completeness does not
imply a primitive-recursive search that finds $h$ from $p,q$ in every presentation.  The certificate
may be an input or the output of a separately justified algorithm.

\subsection{Presented algebras over a certified base ring}

The same construction works over any fixed certified commutative ring $R$.  Polynomial coefficients
are $R$-codes; coefficient identities carry $R$-equality witnesses; and the ideal witness includes
both multiplier polynomials and the coefficientwise equality certificates.  A uniform theorem may
quantify over a finite presentation code only when the base-ring operations and their correctness
proofs are supplied through a fixed interface.  It may not call an arbitrary ring program stored in
the input.

Nested presentations can be flattened.  If
$R=\Z[Y]/J$ and $A=R[X]/I$, replace every coefficient by a $Y$-polynomial and present $A$ directly
as a quotient of $\Z[Y,X]$ by $J$ together with lifted generators of $I$.  The flattening algorithm
is primitive recursive and preserves every supplied ideal witness by substitution.

\subsection{Morphisms and their certificates}

Let
\[
 A=\Z[X_0,\ldots,X_{n-1}]/(g_0,\ldots,g_{r-1}),\qquad
 B=\Z[Y_0,\ldots,Y_{m-1}]/(h_0,\ldots,h_{s-1}).
\]
A presented ring map $\phi:A\to B$ consists of polynomials
$p_0(Y),\ldots,p_{n-1}(Y)$ and, for each $g_i$, a witness
\[
  g_i(p_0,\ldots,p_{n-1})=\sum_{j<s} a_{ij}h_j.
\]
The induced map sends $q(X)$ to $q(p(Y))$.  If $q-q'=\sum b_i g_i$, then substitution gives
\[
 q(p)-q'(p)=\sum_i b_i(p)g_i(p)
           =\sum_j\left(\sum_i b_i(p)a_{ij}\right)h_j,
\]
so the map has an explicit equality-witness transformer.

Composition substitutes the image tuple of one map into the other and composes the relation
certificates by the displayed formula.  Identity maps use the coordinate variables.  An isomorphism
certificate is a pair of presented maps plus coordinatewise witnesses that both composites equal the
identities.  Thus affine morphisms and affine isomorphisms are finite arithmetic data.

\subsection{Modules, matrices, kernels, and cokernels}

A finitely presented $A$-module is given by a matrix
$R:A^u\to A^v$ and the cokernel presentation $A^v/\im R$.  A vector equality witness is a coefficient
vector $z\in A^u$ satisfying $x-y=Rz$, with componentwise ring witnesses.  The same witness algebra
as in Theorem~\ref{thm:quotient-setoid} gives a certified module.

Matrix multiplication, direct sums, tensor-presentation matrices, dual maps, Koszul differentials,
and finite chain complexes are all primitive recursive on finite codes.  Kernel membership is
expressed by a checked equation $Rx=0$; image membership by a witness $x=Ry$; and exactness at a
particular element by an explicit preimage witness.  This is exactly the level needed for finite
\v{C}ech arguments.

\subsection{Gröbner bases as optional certificates}

A Gröbner basis is valuable for producing ideal witnesses and normal forms, but the foundational
interface should not assume that a uniform unbounded Buchberger run is a primitive-recursive
oracle.  Three safe modes are available.

\begin{description}[style=nextline]
\item[Verification mode.]
The input contains a proposed basis, representations of its elements in the original ideal,
reduction traces for the original generators, and reductions of all relevant $S$-polynomials to
zero.  A fixed checker verifies Buchberger's criterion and every identity.

\item[Fuelled mode.]
A fixed Buchberger evaluator runs for a supplied bound and returns its finite trace.  Soundness of a
successful output is proved; failure to finish within the bound makes no completeness claim.

\item[Bounded-family mode.]
For a family with a separately proved degree/size bound, the bound becomes the fuel for a total
primitive-recursive constructor.
\end{description}

This separation matters for computational algebraic geometry.  Ideal membership certificates,
Hilbert-function tables, elimination identities, and finite standard-monomial bases can all be
checked in iPRA even when their general-purpose discovery algorithm is not part of the theorem.

\subsection{A small exact example}

In $\Z[x]/(x^2-2)$ the equality $x^4=4$ is certified by
\[
 x^4-4=(x^2+2)(x^2-2).
\]
The executable companion represents each polynomial sparsely and verifies this coefficientwise.
Nothing is inferred from an approximate root of $x^2-2$; the quotient identity is exact finite
algebra.

\section{Localization as a certified setoid}

\subsection{Fractions and equality witnesses}

Let $A$ be a certified commutative ring and fix $f\in A$.  A code for an element of $A_f$ is a pair
$(a,n)$, written $a/f^n$.  The pair is valid when $a$ is a valid $A$-code.  Equality is not defined by
normalizing fractions.  A witness from $a/f^n$ to $b/f^m$ is
\begin{enumerate}[label=(\alph*)]
\item a natural number $k$, and
\item an $A$-equality witness for
\begin{equation}\label{eq:loc-equality}
             f^k\bigl(f^m a-f^n b\bigr)=0.
\end{equation}
\end{enumerate}
This is the standard equality relation for localization at
$\{1,f,f^2,\ldots\}$, made proof-relevant.

\begin{theorem}[Localization witness calculus]\label{thm:localization-calculus}
The relation \eqref{eq:loc-equality} is an equivalence relation.  Addition and multiplication
\[
 \frac a{f^n}+\frac b{f^m}
   =\frac{f^m a+f^n b}{f^{n+m}},\qquad
 \frac a{f^n}\frac b{f^m}=\frac{ab}{f^{n+m}}
\]
are certified congruent operations.  All witness transformers are primitive recursive.
\end{theorem}

\begin{proof}
For reflexivity choose $k=0$ and the certificate for $f^n a-f^n a=0$.  Symmetry uses the same
$k$ and negates the base equality.

For transitivity, suppose
\begin{align*}
 f^k(f^m a-f^n b)&=0,\\
 f^\ell(f^r b-f^m c)&=0.
\end{align*}
Multiply the first identity by $f^{\ell+r}$ and the second by $f^{k+n}$, then add.  The middle
terms cancel and give
\[
 f^{k+\ell+m}(f^r a-f^n c)=0.
\]
Thus a valid killing exponent for the composite witness is $k+\ell+m$, where $m$ is the
denominator exponent of the middle fraction.  The new base-ring witness is constructed from the old
ones using multiplication congruence, addition, and the distributive-law certificates.

For addition and multiplication, expand the cross-multiplied difference of the outputs.  It is a
finite sum of multiples of the two input differences.  If input witnesses kill those differences by
$f^k$ and $f^\ell$, then multiplying by a sufficiently large explicit power, for example
$f^{k+\ell+n+m+n'+m'}$, kills the output difference.  A sharper exponent is easy but unnecessary.
The expansion and witness combination are fixed polynomial identities, hence primitive recursive.
\end{proof}

The theorem remains valid with zero divisors.  The killing exponent is essential precisely because
cross multiplication need not be cancellative.

\subsection{Localization of modules}

For an $A$-module $M$, represent an element of $M_f$ by $x/f^n$ and witness equality by
\[
             f^k(f^m x-f^n y)=0_M.
\]
The preceding proof uses only additivity and scalar multiplication, so it carries over verbatim.
The restriction map $M\to M_f$ sends $x$ to $x/f^0$.  Equality of two restricted elements is
witnessed by a local annihilation equation $f^k(x-y)=0$.

If $g\in A$, restriction from $M_f$ to $M_{fg}$ sends
\[
             \frac{x}{f^n}\longmapsto \frac{g^n x}{(fg)^n}.
\]
This formula is preferable to an abstract universal property in an object-theory implementation:
it is a fixed primitive-recursive map, and its functoriality is a finite polynomial identity.

\subsection{Universal property in certificate form}

Suppose $\phi:A\to B$ is a certified ring map and $u$ is a supplied inverse witness for $\phi(f)$,
that is, a code $v\in B$ with certificates
\[
          \phi(f)v=1,
          \qquad v\phi(f)=1.
\]
Then
\[
       \widetilde\phi(a/f^n)=\phi(a)v^n
\]
defines a certified map $A_f\to B$.  If \eqref{eq:loc-equality} holds, applying $\phi$ and multiplying
by the appropriate power of $v$ proves equality of the images.  Uniqueness is checked on a supplied
fraction by the equation
\[
 a/f^n=(a/1)(1/f)^n.
\]
Thus the universal property can be used in iPRA as a family of finite witness transformations,
without quantifying over a completed class of all ring maps.

\subsection{Localizing a presented quotient}

Let $A=R[X]/I$ be represented by ideal-combination witnesses.  A localization equality witness for
$p/f^n$ and $q/f^m$ consists of a killing exponent $k$ and multipliers $h_i$ satisfying
\[
 f^k(f^m p-f^n q)=\sum_i h_i g_i
\]
in the polynomial ring.  This is a single coefficientwise certificate.  Hence quotient formation and
localization compose without choosing normal forms.

The same presentation can be flattened algebraically:
\[
        (R[X]/I)_f \cong R[X,T]/(I,\,Tf-1).
\]
The map from a fraction sends $p/f^n$ to $pT^n$; the inverse sends $T$ to $1/f$.  Both composites have
explicit relation witnesses.  The fraction presentation is usually smaller for gluing, while the
polynomial presentation is convenient for elimination and finite-presentation libraries.

\subsection{Primitive-recursive implementation and induction audit}

The code operations needed here are bounded exponentiation, tuple access, polynomial addition and
multiplication, and witness-vector combination.  Exponentiation is primitive recursion on $n$ with
the quantifier-free invariant that the accumulator equals the requested finite power.  Polynomial
operations are bounded folds over finite lists.  No induction formula asserts that two localization
classes are equal; it asserts only that a concrete, checked equality witness has been transformed
correctly.

\section{The explicit power-cover lemma}

\subsection{Unit-ideal covers}

For a ring $A$, a finite family of principal opens $D(f_0),\ldots,D(f_{r-1})$ covers
$\Spec(A)$ exactly when the elements $f_i$ generate the unit ideal.  In the present framework the
cover is not represented by quantifying over prime ideals.  It is represented by coefficients
$c_i\in A$ and an equality witness for
\begin{equation}\label{eq:unit-cover}
                 \sum_{i<r} c_i f_i=1.
\end{equation}
This certificate is finite, stable under ring maps, and directly usable by an arithmetic proof.

Gluing fractions usually requires a common denominator power.  The same principal opens are also a
cover after replacing every $f_i$ by $f_i^N$, but an iPRA proof needs explicit coefficients.

\begin{theorem}[Power-cover constructor]\label{thm:power-cover}
Let $r\geq1$, let $f_i,c_i\in A$ satisfy \eqref{eq:unit-cover}, and let $N\geq1$.  There is an
explicit primitive-recursive construction of coefficients $d_i$ such that
\begin{equation}\label{eq:power-cover}
                  \sum_{i<r}d_i f_i^N=1.
\end{equation}
The construction is uniform in the finite ring codes, $r$, and $N$.
\end{theorem}

\begin{proof}
Set
\[
                    M=r(N-1)+1.
\]
Expand the finite power of \eqref{eq:unit-cover}:
\begin{equation}\label{eq:multinomial-cover}
  1=\left(\sum_{i<r}c_i f_i\right)^M
   =\sum_{\substack{\alpha\in\N^r\\ |\alpha|=M}}
      \binom{M}{\alpha}\prod_{j<r}c_j^{\alpha_j}f_j^{\alpha_j}.
\end{equation}
For every exponent vector $\alpha$ of total degree $M$, some coordinate satisfies
$\alpha_i\geq N$.  Otherwise
$|\alpha|\leq r(N-1)<M$.  Choose the least such $i$ and assign the corresponding term of
\eqref{eq:multinomial-cover} to bucket $i$.  Factor $f_i^N$ from every term in bucket $i$ and let
$d_i$ be the sum of the residual factors.  Then the buckets partition the finite expansion, so
\eqref{eq:power-cover} follows.
\end{proof}

The proof is constructive even over rings with zero divisors.  It uses no radical membership,
prime ideal, compactness, or excluded middle beyond comparisons of natural-number exponents.

\subsection{Exact coefficient formula}

For reproducibility, define
\[
  \iota(\alpha)=\min\set{i<r:\alpha_i\geq N}.
\]
Then
\begin{equation}\label{eq:di-formula}
 d_i=\sum_{\substack{|\alpha|=M\\ \iota(\alpha)=i}}
 \binom{M}{\alpha}
 c_i^{\alpha_i}f_i^{\alpha_i-N}
 \prod_{j\neq i}c_j^{\alpha_j}f_j^{\alpha_j}.
\end{equation}
The finite weak compositions of $M$ into $r$ parts can be enumerated by nested primitive recursion.
The multinomial coefficient is computed from bounded products and exact divisions, or avoided
entirely by an iterative expansion that carries multiplicities by repeated addition.

\subsection{An iPRA-friendly iterative proof}

Formula \eqref{eq:di-formula} is mathematically transparent, but the object proof need not formalize
factorials.  A more direct implementation maintains a finite polynomial table indexed by exponent
vectors with coordinates at most $M$.  Starting from the constant table $1$, perform $M$ rounds of
multiplication by $\sum c_i f_i$.  The quantifier-free invariant states:

\begin{enumerate}[label=(\roman*)]
\item every occupied exponent vector has total degree equal to the round number;
\item the sum of table entries times their monomials equals the current power;
\item each entry is the accumulated coefficient obtained by the fixed update rule.
\end{enumerate}

After round $M$, scan the finite table, choose the least coordinate at least $N$, and add the factored
term to the corresponding $d_i$.  A second quantifier-free invariant states that the processed prefix
of the table equals the partial sum $\sum_i d_i f_i^N$.  This route uses only finite arrays, bounded
comparison, and ring operations already available in the presentation.

\subsection{Examples and validation}

For $f_0=2$, $f_1=3$, $c_0=-1$, and $c_1=1$, the basic cover certificate is
$-2+3=1$.  At power $N=2$, the constructor used in the executable companion returns
\[
                   d_0=7,\qquad d_1=-3,
\]
and indeed
\[
                    7\cdot2^2-3\cdot3^2=28-27=1.
\]

The same algorithm was tested in $\Z/12\Z$ with $f_0=3$, $f_1=4$ and
$c_0=3$, $c_1=1$, since $3\cdot3+1\cdot4=1\pmod {12}$.  This test is useful because neither
principal element is cancellable and localization witnesses genuinely need annihilating powers.

\subsection{Functoriality}

If $\phi:A\to B$ is a certified ring map, applying $\phi$ to the input cover and to every operation
in \eqref{eq:di-formula} gives exactly the power-cover coefficients constructed in $B$.  Thus the
constructor is functorial up to the supplied equality witnesses.  This fact later makes atlas
refinement and base change finite certificate transformations rather than new existence proofs.

\section{Explicit principal-cover gluing}

\subsection{Statement}

Let $A$ be a commutative ring, $M$ an $A$-module, and suppose
\[
                  \sum_{i<r}c_i f_i=1.
\]
For each $i$, let a local section on $D(f_i)$ be represented by
\[
                  s_i=\frac{m_i}{f_i^{n_i}}\in M_{f_i}.
\]
On the overlap $D(f_if_j)$, the two restrictions are represented after denominator
homogenization by fractions in the single localization $M_{f_if_j}$.

\begin{theorem}[Primitive-recursive Zariski gluing]\label{thm:principal-gluing}
Assume that for every $i,j<r$ a finite localization witness proves that $s_i$ and $s_j$ agree in
$M_{f_if_j}$.  Then an explicit primitive-recursive formula produces $m\in M$ and witnesses that
$m/1=s_i$ in every $M_{f_i}$.  If two global sections have the same restrictions, an explicit
primitive-recursive witness proves them equal.  The theorem is valid over rings with zero divisors.
\end{theorem}

This is the sheaf equalizer theorem for the quasi-coherent module associated to $M$, presented in a
form directly usable by iPRA.

\subsection{Denominator homogenization}

Set
\[
                         N=\max_{i<r} n_i,
          \qquad a_i=f_i^{N-n_i}m_i.
\]
Then $s_i=a_i/f_i^N$.  An overlap witness can be written as an exponent $k_{ij}$ satisfying
\begin{equation}\label{eq:overlap-kill}
  (f_if_j)^{k_{ij}}
  \bigl(f_j^N a_i-f_i^N a_j\bigr)=0_M.
\end{equation}
The witness includes the module-equality certificate for this equation.  Let
\[
                         K=\max_{i,j<r}k_{ij},
             \qquad q=N+K+1,
             \qquad b_i=f_i^{q-N}a_i=f_i^{K+1}a_i.
\]
Multiplying \eqref{eq:overlap-kill} by the remaining powers needed to reach $K$ gives
\begin{equation}\label{eq:cross-common}
                         f_j^q b_i=f_i^q b_j
                         \qquad(i,j<r).
\end{equation}
Indeed the difference of the two sides equals
\[
 f_if_j(f_if_j)^K(f_j^Na_i-f_i^Na_j),
\]
which is zero by the synchronized overlap witness.  The harmless extra $+1$ in $q$ avoids a special
case when every denominator and every killing exponent is zero.

\subsection{Construction of the global section}

Apply Theorem~\ref{thm:power-cover} to the exponent $q$.  It returns $d_i$ satisfying
\[
                         \sum_{i<r}d_i f_i^q=1.
\]
Define
\begin{equation}\label{eq:global-glue}
                         m=\sum_{i<r}d_i b_i.
\end{equation}
For a fixed chart $j$, use \eqref{eq:cross-common}:
\begin{align*}
 f_j^q m
   &=\sum_i d_i f_j^q b_i\\
   &=\sum_i d_i f_i^q b_j\\
   &=b_j.
\end{align*}
Thus $m/1=b_j/f_j^q=a_j/f_j^N=m_j/f_j^{n_j}$ in $M_{f_j}$.  Every equality in this calculation is
accompanied by a finite witness obtained from module distributivity, the pairwise cross witnesses,
and the power-cover certificate.

This proves existence without search.  Once the overlap witnesses and the unit-ideal cover are in
hand, the global section is a fixed bounded formula.

\subsection{Uniqueness}

Suppose $m,m'\in M$ have equal restrictions in every $M_{f_i}$.  For each $i$, let $e_i$ witness
\[
                            f_i^{e_i}(m-m')=0.
\]
Set $E=1+\max_i e_i$.  Then $f_i^E(m-m')=0$ for every $i$.  Use the power-cover constructor to find
$u_i$ with $\sum u_i f_i^E=1$.  Therefore
\[
  m-m'=\sum_i u_i f_i^E(m-m')=0.
\]
This is a direct equality witness in $M$.  No cancellation, prime ideal argument, or excluded middle
is used.

\begin{corollary}[Finite principal-cover equalizer]\label{cor:equalizer}
The sequence
\[
  0\longrightarrow M
   \longrightarrow \prod_{i<r}M_{f_i}
   \rightrightarrows \prod_{i,j<r}M_{f_if_j}
\]
is exact in the following certificate sense:
\begin{enumerate}[label=(\alph*)]
\item the first map has an explicit joint-monicity witness transformer;
\item every tuple with supplied pairwise compatibility witnesses has an explicitly constructed
      preimage;
\item the restriction witnesses and uniqueness witness are primitive recursive in the input codes.
\end{enumerate}
\end{corollary}

Finite products are coded by lists, so the theorem is a statement entirely about natural numbers.
The use of the word ``exact'' does not import an abstract abelian-category library.

\subsection{A zero-divisor example}

Take $A=M=\Z/12\Z$, $f_0=3$, $f_1=4$, and cover coefficients $c_0=3$, $c_1=1$:
\[
                    3\cdot3+1\cdot4=13=1\pmod {12}.
\]
The local numerators $9$ on $D(3)$ and $8$ on $D(4)$ both represent the global class $5$:
\[
                  3(9-5)=0,\qquad 4(8-5)=0\pmod {12}.
\]
On the overlap, $(3\cdot4)^1=0$ kills every difference, so the compatibility witness has exponent
one.  The gluing algorithm takes $N=0$, $K=1$, $q=2$ if the displayed formula is used literally.
A slightly optimized implementation uses $q=1$ because the synchronized cross identity already
holds; either choice is sound.  The executable companion uses the optimized exponent and returns the
unique global class $5$.

This example demonstrates why localization equality must retain killing exponents.  Cross
multiplication alone would not be a cancellative equality test in a ring with zero divisors.

\subsection{Primitive-recursive complexity}

All loops are bounded by primitive-recursive terms in the input lengths and exponents:

\begin{itemize}
\item scanning the $n_i$ and $k_{ij}$ uses bounded maxima;
\item exponentiation uses primitive recursion;
\item the power-cover constructor enumerates finitely many weak compositions or performs a bounded
      iterative expansion;
\item the final sums are finite module folds.
\end{itemize}

A convenient quantifier-free invariant for the final fold over $i<t$ is
\[
 f_j^q m_t=\left(\sum_{i<t}d_i f_i^q\right)b_j,
\]
where $m_t=\sum_{i<t}d_i b_i$.  The invariant is an equality of coded module terms with a concrete
witness accumulator.  It contains no quantification over local sections or opens.

\subsection{Refinements and iterated covers}

If a chart $D(f_i)$ is further covered by principal opens $D(f_i g_{i\alpha})$ with a local
unit-ideal certificate in $A_{f_i}$, clear the finitely many denominators in that certificate and
obtain a power cover in $A$.  Gluing first within each chart and then across the outer cover agrees
with gluing on the flattened refinement.  The proof follows uniqueness: both constructions have the
same restrictions on every refined principal piece, so the global outputs are equal.

Thus finite cover refinement is coherent without choosing canonical normal forms.  This is the
basic associativity needed to build finite affine atlases and sheaf data.

\section{Certified finite affine atlases}

\subsection{Strict principal atlases}

The simplest geometric gluing object has one principal overlap for each ordered pair of charts.

\begin{definition}[Strict certified affine atlas]\label{def:strict-atlas}
A strict certified affine atlas consists of:
\begin{enumerate}[label=(\alph*)]
\item a positive integer $r$ and finitely presented rings $A_i$ for $i<r$;
\item elements $f_{ij}\in A_i$ for every $i,j<r$, with $f_{ii}=1$;
\item certified isomorphisms
\[
           \phi_{ij}:(A_i)_{f_{ij}}\xrightarrow{\sim}(A_j)_{f_{ji}};
\]
\item identity witnesses for $\phi_{ii}$ and inverse witnesses for
      $\phi_{ji}\phi_{ij}=\mathrm{id}$;
\item on every triple $(i,j,k)$, a finite common localization and a witness that
\[
                         \phi_{ik}=\phi_{jk}\phi_{ij}
\]
there.
\end{enumerate}
Every ring, element, map, localization, and equality in this definition is a natural-number code
with the certificates described in the preceding sections.
\end{definition}

Externally, the usual gluing theorem for schemes turns this data into a scheme covered by the affine
schemes $\Spec(A_i)$.  Internally, no topological quotient is needed.  A theorem about the glued
object is translated into a finite theorem about the charts and the certified overlap maps.

The strict form already covers projective space: on the chart $X_i\neq0$, the overlap with the chart
$X_j\neq0$ is the principal open where $X_j/X_i$ is invertible.

\subsection{Finite principal-overlap atlases}

Pairwise intersections of affine opens need not always be one principal open.  The practical
extension replaces $f_{ij}$ by a finite family $f_{ij\alpha}$, each giving a principal overlap piece.
The data include:

\begin{itemize}
\item a unit-ideal or relative-cover certificate that the pieces cover the intended intersection;
\item a localization isomorphism on every piece;
\item compatibility witnesses on double refinements;
\item cocycle witnesses on triple refinements.
\end{itemize}

The principal-cover gluing theorem proves that chartwise constructions independent of the chosen
piece descend to the whole intersection.  Refinement coherence follows from uniqueness.  Thus the
additional indexing is finite bookkeeping rather than a new logical principle.

\subsection{The internal geometric language}

A finite atlas supports the following notions without constructing an internal prime spectrum.

\begin{description}[style=nextline]
\item[Regular function.]
A tuple $a_i\in A_i$ with localization equality witnesses
$\phi_{ij}(a_i)=a_j$ on overlaps.  The principal-cover theorem glues such tuples whenever one chart
is itself presented by a finite principal cover.

\item[Morphism.]
A chartwise finite family of presented ring maps in the contravariant direction, together with a
finite target-chart cover and overlap commutativity witnesses.  Relational finite-use composition
may choose among several valid target charts; the induced map on an invariant observation is
choice-free.

\item[Open subobject.]
A finite family of principal elements on each chart with compatibility certificates under overlap
maps.  Union and intersection are computed by concatenation and products of principal elements.

\item[Closed subobject.]
A finite ideal presentation on each chart, stable under localization and equipped with overlap
identities.  A global finitely presented closed subscheme is represented by finitely many local
relation lists.

\item[Point over a presented ring $R$.]
A compatible family of certified ring maps $A_i\to R$ on a selected chart cover, or more commonly a
single chart map with a certificate that its denominators are invertible.  Projective homogeneous
coordinates are a compact specialized presentation of the same data.
\end{description}

These are theorem-relative representations.  A statement requiring only chartwise equations never
needs to enumerate the topological points of the scheme.

\subsection{Atlas morphisms and refinement}

A refinement replaces a chart $A_i$ by a finite principal cover.  Its certificate contains the
unit-ideal relation and the localization maps.  A common refinement of two finite atlases is formed
by taking finite pairwise intersections and then principal refinements.  All index bounds are
primitive recursive in the atlas sizes.

\begin{proposition}[Refinement invariance]\label{prop:refinement-invariance}
Suppose a finite atlas observation is defined by compatible local data and has uniqueness under the
principal-cover equalizer.  Then its value is invariant, with an explicit equality witness, under
finite principal refinement.
\end{proposition}

\begin{proof}
Pull the data to the common refinement.  Both candidate values have equal restrictions on every
refined principal piece.  Apply principal-cover uniqueness chart by chart and then on the outer
finite atlas.  The resulting equality witness is obtained by nested applications of the power-cover
constructor.
\end{proof}

This proposition supplies the practical replacement for proving equality of two constructions on
all open sets.

\subsection{Base change and fiber products}

For presented affine rings, tensor products have explicit presentations.  If
$A\leftarrow R\to B$ are finitely presented, then $A\otimes_R B$ is presented by the combined
variables and relations, together with equations identifying the two images of the generators of
$R$.  The construction is primitive recursive on the finite presentations.

Base-changing an atlas along a presented morphism applies this tensor construction chartwise and
base-changes each localization certificate.  Fiber products of finite atlases are obtained by the
finite product of chart index sets.  Overlap cocycles are inherited by functoriality of localization
and tensor-product maps.  Thus many finite-limit constructions in algebraic geometry fit the same
certificate language.

\subsection{Properties as finite certificates}

Some geometric properties have direct finite witnesses on a chosen atlas:

\begin{itemize}
\item \emph{finite presentation}: already part of every chart and morphism code;
\item \emph{separatedness}: a finite presentation/certificate for each affine diagonal or,
      in common cases, affine pairwise intersections with multiplication-map surjectivity witnesses;
\item \emph{smoothness of a fixed relative dimension}: a finite principal cover on which specified
      Jacobian minors are invertible, together with identities certifying the relations and rank;
\item \emph{finite locally free module of rank $d$}: local bases and inverse matrices on a finite
      principal cover;
\item \emph{proper projective presentation}: a closed homogeneous presentation inside a fixed
      projective space.
\end{itemize}

These certificates are sufficient conditions with soundness proofs.  The framework does not assert
that a uniform primitive-recursive algorithm discovers such certificates for every input
presentation.

\subsection{Conventional semantics}

Externally, Definition~\ref{def:strict-atlas} is interpreted by the standard gluing theorem for
schemes.  The finite ring maps determine open immersions of affine schemes, and the inverse and
cocycle witnesses provide the usual gluing datum.  Every internal chartwise theorem is sound for the
resulting conventional scheme because its algebraic restriction maps are exactly the localization
maps used by the structure sheaf.

A full formal equivalence between the finite-atlas category and an appropriate category of
quasi-compact schemes with chosen atlases would be valuable, but it is not required for individual
iPRA theorems.  Each theorem can state its atlas hypotheses explicitly and prove the desired finite
observation directly.

\section{Quasi-coherent modules, vector bundles, and finite \v{C}ech complexes}

\subsection{Affine module data}

On an affine chart $U_i=\Spec(A_i)$, a quasi-coherent module is represented by an $A_i$-module
$M_i$.  For a finitely presented theory one usually takes a finite presentation matrix, although the
principal-cover gluing theorem only needs a certified module interface.  Restriction to a principal
open $D(f)$ is the localization $(M_i)_f$.

For a strict atlas, descent data for a quasi-coherent module consists of certified isomorphisms
\[
   \theta_{ij}:(M_i)_{f_{ij}}
      \xrightarrow{\sim}
      (M_j)_{f_{ji}}
\]
semilinear with respect to $\phi_{ij}$, together with inverse and triple-cocycle witnesses.  With
finite principal-overlap atlases, the same data are supplied on every overlap piece and checked on
finite refinements.

\begin{definition}[Certified quasi-coherent atlas module]
A certified quasi-coherent atlas module is a finite affine atlas, one certified module on each chart,
and finite localization isomorphism/cocycle certificates on the overlaps.
\end{definition}

This definition is purely finite.  Its external interpretation is the conventional quasi-coherent
sheaf obtained by descent.

\subsection{Sections and the sheaf equalizer}

A global section code is a tuple $(s_i)_{i<r}$ with $s_i\in M_i$ and overlap witnesses
\[
                    \theta_{ij}(s_i)=s_j.
\]
If a chart or overlap is represented by several principal pieces, compatibility on those pieces is
compressed by Theorem~\ref{thm:principal-gluing}.  Equality of global sections is componentwise
certificate equality; refinement invariance follows from Proposition~\ref{prop:refinement-invariance}.

For a principal cover of a single affine chart, Corollary~\ref{cor:equalizer} proves the sheaf
condition directly.  Hence the structure presheaf $D(f)\mapsto A_f$ and module presheaf
$D(f)\mapsto M_f$ are sheaves on the finite standard-open basis in exactly the finite sense used by
an iPRA proof.

\subsection{Vector bundles and transition matrices}

A rank-$d$ vector bundle on a finite atlas is represented by a free module $A_i^d$ on each chart and
invertible matrices
\[
                     G_{ij}\in\operatorname{GL}_d((A_i)_{f_{ij}})
\]
with supplied inverse matrices and cocycle witnesses $G_{ik}=G_{jk}G_{ij}$ on triple overlaps.
All matrix identities are finite coefficient identities in localizations.  A section is a tuple of
vectors $v_i\in A_i^d$ satisfying $v_j=G_{ij}v_i$.

Line bundles are the case $d=1$: transition matrices are units.  The twists $\cO(d)$ on projective
space have transition units given by integer powers of coordinate ratios.  Negative powers live in
the overlap localization and require no extra total inversion operation outside its certified
domain.

\subsection{Coherent modules and finite resolutions}

A coherent module may be represented chartwise by a finite presentation
\[
                       A_i^{u_i}\xrightarrow{R_i}A_i^{v_i}\to M_i\to0.
\]
Overlap isomorphisms are matrices on generators together with witnesses that they respect relations,
plus inverse witnesses.  Kernels, cokernels, tensor products, and duals are represented by explicit
finite matrices and relation certificates.

A finite locally free resolution is especially useful: every differential is a matrix, the equation
$d_{p-1}d_p=0$ is coefficientwise, and exactness at a supplied vector is certified by a preimage.
The framework can verify a finite resolution without deciding module equality globally.

\subsection{Finite \v{C}ech complexes}

Let $U_0,\ldots,U_{r-1}$ be a finite atlas or a finite principal cover.  For a module datum
$\mathcal F$, define the ordered \v{C}ech group in degree $p$ by the finite direct sum
\[
 \check C^p(U,\mathcal F)
   =\bigoplus_{i_0<\cdots<i_p}
     \mathcal F(U_{i_0}\cap\cdots\cap U_{i_p}).
\]
Each multiple intersection is represented by an iterated localization or a finite principal
refinement.  The differential is
\[
   (ds)_{i_0\ldots i_{p+1}}
      =\sum_{a=0}^{p+1}(-1)^a
       s_{i_0\ldots\widehat{i_a}\ldots i_{p+1}}\big|_{U_{i_0\ldots i_{p+1}}}.
\]
All sums and restrictions are fixed PR operations on finite tuples.

The identity $d^2=0$ is proved by pairing the two ways of deleting two indices; their signs are
opposite.  An iPRA proof enumerates the bounded pairs $(a,b)$ and carries a quantifier-free
coefficient accumulator.  It does not use induction on a statement quantifying over sheaves.

\subsection{Cohomology by explicit normal forms}

For a finite complex, a cohomology class may be represented in several certificate-relative ways:

\begin{enumerate}[label=\arabic*.]
\item a cocycle $z$ with a witness $dz=0$, modulo proof-relevant equality
      $z-z'=dy$;
\item a supplied direct-sum normal form with constructors for decomposition and uniqueness;
\item a finite presentation matrix for $\ker d/\im d$;
\item a contracting homotopy $h$ satisfying $hd+dh=1$ on a designated summand.
\end{enumerate}

The projective-space computation uses the fourth and second forms.  Every monomial-weight summand is
finite and has either a one-dimensional survivor or an explicit insertion homotopy.

\subsection{Relation to conventional sheaf cohomology}

The arithmetic construction gives exact \v{C}ech cohomology of a chosen finite atlas.  To identify it
with derived-functor sheaf cohomology one needs a comparison theorem, usually supplied by affine
acyclicity for quasi-coherent sheaves and acyclicity of finite intersections.  For the standard cover
of projective space, this is the standard conventional argument.

There are two honest formalization targets:

\begin{description}[style=nextline]
\item[Immediate target.]
Formalize the finite \v{C}ech complex and its normal forms in iPRA.  This already proves concrete
local-to-global statements and computes a canonical finite algebraic invariant.

\item[Extended target.]
Develop a coded homological-algebra library and an affine-acyclicity theorem, then prove that the
finite \v{C}ech result represents the usual sheaf cohomology.  This is a larger project and is not
silently assumed to be complete here.
\end{description}

The report states conventional cohomology conclusions only with this comparison boundary made
explicit.

\section{Projective schemes and computational algebraic geometry}

\subsection{The standard affine charts of \texorpdfstring{$\Proj$}{Proj}}

Let
\[
 S=R[X_0,\ldots,X_n]/I
\]
be a finitely presented graded ring, with $I$ given by finitely many homogeneous relations.  The
standard chart $D_+(X_i)$ is represented algebraically by the degree-zero part of the localization
$S_{X_i}$.  For a projective subscheme of $\mathbf P^n_R$, this chart has the familiar finite
presentation obtained by setting $X_i=1$ and introducing affine coordinates
\[
                    x_{ij}=X_j/X_i\qquad(j\neq i).
\]
Each homogeneous relation is dehomogenized with respect to $X_i$.

The overlap of charts $i$ and $j$ is the principal open in chart $i$ where $x_{ij}$ is invertible.
The transition map is
\[
 x_{jk}\longmapsto x_{ik}/x_{ij},
 \qquad
 x_{ji}\longmapsto 1/x_{ij}.
\]
These formulas are maps of localization presentations.  Their inverse and triple-cocycle equations
are finite rational-function identities, hence localization equality certificates.

\begin{proposition}[Finite atlas for a projective presentation]\label{prop:proj-atlas}
For every fixed finite homogeneous presentation of $S$, the $n+1$ standard charts, their principal
overlaps, transition maps, inverse witnesses, and triple cocycles are primitive-recursively
constructible.  They form a strict certified affine atlas.
\end{proposition}

\begin{proof}
Dehomogenization is a bounded substitution on each finite polynomial code.  The transition maps are
fixed monomial substitutions into localizations.  Applying a homogeneous relation and clearing the
explicit denominator gives its target relation certificate.  The inverse and cocycle equations are
monomial exponent identities; a bounded coefficient checker verifies them.  The number of charts
and overlaps is finite and bounded by polynomial expressions in $n$.
\end{proof}

Thus a large class of projective geometry can be developed without defining $\Proj$ internally as a
set of homogeneous prime ideals.

\subsection{Homogeneous coordinates and points}

For a field or domain with a suitable exact presentation, a projective point may be represented by a
nonzero tuple $(x_0,\ldots,x_n)$ modulo common scaling.  The existing Diophantine geometry packet
uses primitive representatives and height.  The atlas representation and homogeneous-coordinate
representation are complementary:

\begin{itemize}
\item homogeneous coordinates are compact for point enumeration and height arguments;
\item affine charts are natural for local rings, functions, sheaves, Jacobians, and descent;
\item a point with $x_i$ invertible maps to chart $i$ by the ratios $x_j/x_i$;
\item two chart descriptions agree by the certified transition maps.
\end{itemize}

A theorem may use a finite relational choice of a nonzero coordinate.  If its final output is
projectively invariant, the finite-use framework eliminates dependence on the selected chart.
Alternatively, a primitive representative over $\Z$ can choose the least nonzero coordinate
deterministically.

\subsection{Projective morphisms}

A tuple of homogeneous forms $F_0,\ldots,F_m$ of one common degree defines a projective morphism on
the locus where they do not vanish simultaneously.  A finite certificate for a global morphism on a
projective subscheme consists of:

\begin{enumerate}[label=(\alph*)]
\item relation witnesses showing that the image satisfies the target homogeneous equations;
\item a finite principal-cover certificate that the $F_j$ generate the unit ideal on each source
      chart, or equivalent local nonvanishing certificates;
\item chartwise rational formulas and overlap commutativity witnesses.
\end{enumerate}

Composition substitutes homogeneous forms and multiplies degrees.  The relation and base-locus
certificates transform by finite polynomial identities.  Height upper bounds are obtained from
coefficient norms exactly as in the Diophantine packet; lower bounds still require gcd/resultant or
base-locus information and are not produced by the atlas formalism alone.

\subsection{Jacobian and smoothness certificates}

For a closed affine presentation $A=R[X]/(g_1,\ldots,g_s)$, the Jacobian matrix is a finite matrix of
formal partial derivatives.  A certificate that a chosen $d\times d$ minor is invertible on a
principal open consists of a localization inverse witness.  A finite list of such minors and a
unit-ideal cover certificate proves that the required rank condition holds everywhere on the chosen
atlas.

This gives an iPRA-friendly sufficient interface for smoothness in standard complete-intersection
situations.  It does not assert a uniform algorithm deciding smoothness of arbitrary presentations.
The input may carry a Gröbner/elimination certificate producing the relevant Jacobian ideal, or the
minor cover may simply be part of the theorem hypotheses.

\subsection{Divisors and line bundles}

A Cartier divisor on a finite atlas is represented by local nonzerodivisor fractions $g_i$ and unit
transition witnesses $g_i/g_j\in\cO^\times(U_i\cap U_j)$.  The associated line bundle uses those
units as rank-one transition matrices.  Tensor product adds divisor data and multiplies transition
units; dualization inverts them on overlaps.

For projective space, the hyperplane divisor gives $\cO(1)$, and $\cO(d)$ has transition factor
$(X_j/X_i)^d$.  All integer powers are finite localization terms.  The projective-line and
projective-space computations below therefore require no abstract Picard group.

\subsection{Hilbert functions and finite graded certificates}

A homogeneous Gröbner certificate can provide a finite monomial ideal of leading terms.  For a fixed
degree $D$, the standard monomials of degree $D$ are a bounded finite set, and coefficient reduction
proves a spanning or basis statement for the degree-$D$ quotient.  Hence the Hilbert function at a
specified degree is a finite iPRA observation.

A claim about the eventual Hilbert polynomial needs an additional stabilization certificate: a
bound $D_0$, a polynomial $P$, and finite algebraic data proving the relevant recurrence or graded
resolution for all $D\geq D_0$.  The framework checks such a certificate; it does not infer an
unbounded stabilization theorem from a finite table.

\subsection{Elimination, images, and fibers}

For a presented morphism, an elimination certificate can describe equations for a scheme-theoretic
image or a fiber.  Soundness is again finite: substitute the morphism tuple and verify ideal
membership.  Completeness of the proposed elimination ideal requires a Gröbner basis, resultant, or
other finite certificate appropriate to the family.

This division of labor is useful in iPRA:

\begin{itemize}
\item polynomial substitution and certificate verification are generic;
\item the domain-specific algorithm produces the elimination certificate;
\item the geometric theorem consumes only the certified presentation of the resulting image or
      fiber.
\end{itemize}

\subsection{What part of algebraic geometry is thereby accessible?}

The finite-atlas language directly supports substantial fragments of:

\begin{itemize}
\item affine and projective varieties/schemes of finite presentation;
\item regular and rational maps on certified domains;
\item finite morphisms and finite module algebras;
\item closed immersions, fiber products, and base change;
\item vector bundles, line bundles, and finite \v{C}ech complexes;
\item Jacobian calculations and chartwise smoothness certificates;
\item exact finite-field point counting by bounded enumeration;
\item rational/integral point searches once an external height bound is supplied;
\item graded pieces, Hilbert-function observations, and certificate-relative elimination.
\end{itemize}

The framework is intentionally less ambitious about arbitrary non-quasi-compact schemes, unbounded
sites, derived categories as completed universes, and existence theorems whose proofs provide no
primitive-recursive quantitative data.  Those topics can still have finite theorem-specific
observations, but they are not forced into a single universal code.

\section{Case study: \texorpdfstring{$\mathbf P^1$}{P1} and every twist \texorpdfstring{$\cO(d)$}{O(d)}}

\subsection{The two-chart atlas}

Fix a commutative ring $R$ with a certified presentation.  The projective line has two standard
charts
\[
   U_0=\Spec R[t],\qquad U_1=\Spec R[u],
\]
with overlap
\[
   U_{01}=\Spec R[t,t^{-1}]=\Spec R[u,u^{-1}],
   \qquad u=t^{-1}.
\]
The overlap maps are localizations at $t$ and $u$, followed by the certified isomorphism sending
$u$ to $t^{-1}$.  The inverse sends $t$ to $u^{-1}$.  Their composite identities are literal
Laurent-monomial identities.

The line bundle $\cO(d)$ is trivial on each chart.  Let $e_0$ and $e_1$ be the chosen local bases.
On the overlap,
\[
                         e_1=t^d e_0.
\]
This formula is valid for every $d\in\Z$ because $t$ is invertible on the overlap.

A degree-zero \v{C}ech cochain is a pair
\[
                   a(t)e_0\in R[t]e_0,
              \qquad b(u)e_1\in R[u]e_1.
\]
After writing $u=t^{-1}$, the differential is
\begin{equation}\label{eq:p1-delta}
  \delta_d(a,b)=a(t)-t^d b(t^{-1})
        \in R[t,t^{-1}]e_0.
\end{equation}
Since the cover has two members, the ordered \v{C}ech complex is
\[
  0\longrightarrow R[t]\oplus R[u]
    \xrightarrow{\delta_d} R[t,t^{-1}]
    \longrightarrow0.
\]

\subsection{Global sections}

\begin{theorem}[Global sections of $\cO(d)$ on $\mathbf P^1$]\label{thm:p1-h0}
The kernel of $\delta_d$ has the following explicit normal form:
\[
 \check H^0(\mathbf P^1_R,\cO(d))\cong
 \begin{cases}
   R^{d+1},& d\geq0,\\
   0,&d<0.
 \end{cases}
\]
For $d\geq0$, the vector $(c_0,\ldots,c_d)$ corresponds to
\begin{align*}
 a(t)&=\sum_{e=0}^d c_e t^e,\\
 b(u)&=\sum_{j=0}^d c_{d-j}u^j.
\end{align*}
The conversion in both directions is primitive recursive on finite polynomial codes.
\end{theorem}

\begin{proof}
The equality $\delta_d(a,b)=0$ says
\[
                         a(t)=t^d b(t^{-1}).
\]
The left side has Laurent support in $\{e:e\geq0\}$; the right side has support in
$\{e:e\leq d\}$.  Equality of Laurent polynomials is coefficientwise, so the common support lies in
\[
                         \{e:0\leq e\leq d\}.
\]
If $d<0$, this set is empty and both polynomials are zero.  If $d\geq0$, the common Laurent
polynomial has unique coefficients $c_0,\ldots,c_d$.  Rewriting its coefficient at exponent $e$ as
the coefficient of $u^{d-e}$ gives the displayed formula for $b$.

No cancellation in $R$ is used; the conclusion follows from the free Laurent-monomial
presentation.  Extraction and reconstruction scan the finite support interval.
\end{proof}

\begin{corollary}[Global regular functions]\label{cor:p1-global-functions}
For every certified commutative ring $R$,
\[
                         \Gamma(\mathbf P^1_R,\cO)\cong R
\]
in the finite-atlas semantics.  The isomorphism sends a constant $c$ to the pair $(c,c)$.
\end{corollary}

This is already a genuinely scheme-theoretic local-to-global theorem: it is not a statement about a
bounded set of rational points, but about compatible regular functions on two affine charts.

\subsection{First cohomology: the middle-support normal form}

For a Laurent polynomial
\[
                         c(t)=\sum_{e=L}^U c_e t^e,
\]
define three finite pieces:
\begin{align*}
 c_{\geq0}(t)&=\sum_{e\geq0}c_e t^e,\\
 c_{\leq d}(t)&=\sum_{e\leq d}c_e t^e,\\
 c_{(d,0)}(t)&=\sum_{d<e<0}c_e t^e.
\end{align*}
Set
\[
 a(t)=c_{\geq0}(t),
 \qquad
 b(u)=-\sum_{e\leq d}c_e u^{d-e}.
\]
Then
\begin{equation}\label{eq:p1-normal-form}
                     c(t)=\delta_d(a,b)+c_{(d,0)}(t).
\end{equation}

\begin{theorem}[Primitive-recursive $H^1$ normal form]\label{thm:p1-h1}
Every degree-one \v{C}ech cochain has the decomposition
\eqref{eq:p1-normal-form}.  The residual term is supported exactly in the finite interval
$d<e<0$, and two such residual terms differ by a coboundary only when they are coefficientwise equal.
Consequently
\[
 \check H^1(\mathbf P^1_R,\cO(d))\cong
 \begin{cases}
   0,&d\geq-1,\\
   R^{-d-1},&d\leq-2,
 \end{cases}
\]
with basis represented by $t^{d+1},\ldots,t^{-1}$ in the second case.
\end{theorem}

\begin{proof}
The high-support part $c_{\geq0}$ is the $a$ contribution to \eqref{eq:p1-delta}.  For an exponent
$e\leq d$, choose the coefficient of $u^{d-e}$ in $b$ to be $-c_e$.  Then the term
$-t^d b(t^{-1})$ contributes $c_e t^e$.  This proves existence of
\eqref{eq:p1-normal-form}.

The image of the $a$ component is supported in exponents $e\geq0$.  The image of the $b$ component is
supported in exponents $e\leq d$.  Therefore a coboundary has zero coefficient at every exponent
$d<e<0$.  A residual supported in that interval can be a coboundary only when all its coefficients
are zero.  The interval is empty exactly when $d\geq-1$, and otherwise contains
$-d-1$ integers.
\end{proof}

The proof gives algorithms for the quotient map, a section selecting the unique normal form, and
witnesses that their composites are identities.  It is therefore stronger than a mere rank count.

\subsection{Examples}

For $d=2$, a global section is
\[
  a(t)=c_0+c_1t+c_2t^2,
  \qquad
  b(u)=c_2+c_1u+c_0u^2.
\]
There is no $H^1$ obstruction because the interval $2<e<0$ is empty.

For $d=-3$, global sections vanish.  Every Laurent cocycle is cohomologous to a unique expression
\[
                       \alpha t^{-2}+\beta t^{-1},
\]
so $\check H^1\cong R^2$.

The executable companion generated and checked 1,200 random Laurent-polynomial decompositions for
$-7\leq d\leq7$.  The tests use exact integer coefficients and verify reconstruction and support
bounds coefficientwise.

\subsection{iPRA implementation}

A Laurent polynomial code stores a lower exponent, an upper exponent, and a finite coefficient
array, or equivalently a sorted sparse list.  Negative exponents are encoded as signed integers.
The normal-form constructor scans the finite support once and writes each coefficient to one of
three output lists according to the decidable comparisons $e\geq0$ and $e\leq d$.

A quantifier-free invariant after processing a support prefix states:
\begin{enumerate}[label=(\roman*)]
\item the input prefix equals the sum of the partial coboundary and partial residual;
\item the chart-zero list contains only nonnegative exponents;
\item the chart-one list contains only nonnegative $u$-exponents;
\item the residual contains only exponents $d<e<0$.
\end{enumerate}
At the end, coefficientwise uniqueness is proved by a bounded scan of the common finite support
interval.  No induction formula quantifies over polynomials or cohomology classes.

\subsection{From \v{C}ech groups to conventional cohomology}

The standard two-chart cover and its intersection are affine, and $\cO(d)$ is quasi-coherent.  The
usual affine-acyclicity comparison identifies the computed \v{C}ech groups with conventional sheaf
cohomology.  The finite arithmetic proof above is independent of that abstract comparison: it
computes the exact complex and its normal forms.  A full internal development of derived-functor
cohomology would add the comparison theorem as a later layer, not alter the finite calculation.

\section{The standard \v{C}ech computation on \texorpdfstring{$\mathbf P^n$}{Pn}}

\subsection{The finite complex}

Let $n\geq1$ and let
\[
           \mathbf P^n_R=\bigcup_{i=0}^{n}U_i,
           \qquad U_i=D_+(X_i).
\]
Every finite intersection $U_{i_0}\cap\cdots\cap U_{i_p}$ is a standard affine open.  In homogeneous
coordinates, its degree-$d$ sections of $\cO(d)$ are finite Laurent polynomials in which
$X_{i_0},\ldots,X_{i_p}$ may have negative exponents and all other variables have nonnegative
exponents, with total degree $d$.

The ordered \v{C}ech complex is
\begin{equation}\label{eq:pn-cech}
 \check C^p_d=\bigoplus_{0\leq i_0<\cdots<i_p\leq n}
 \left(R[X_0,\ldots,X_n,(X_{i_0}\cdots X_{i_p})^{-1}]\right)_d,
\end{equation}
with the alternating restriction differential.

Every individual code in \eqref{eq:pn-cech} has finite Laurent support.  Thus an iPRA theorem does
not form an actually infinite direct sum: it scans the finitely many exponent vectors occurring in
the input cochain.

\subsection{Decomposition by exponent vector}

Fix $e=(e_0,\ldots,e_n)\in\Z^{n+1}$ with
\[
                            \sum_{i=0}^{n}e_i=d.
\]
Let
\[
                     \operatorname{Neg}(e)=\set{i:e_i<0}.
\]
The monomial $X^e=X_0^{e_0}\cdots X_n^{e_n}$ occurs in the component indexed by
$I=\{i_0<\cdots<i_p\}$ exactly when
\[
                       \operatorname{Neg}(e)\subseteq I.
\]
Consequently the full complex is the coefficientwise direct sum of finite complexes
$C^\bullet(e)$, each having one copy of $R$ for every subset $I$ that contains
$\operatorname{Neg}(e)$.

The differential on $C^\bullet(e)$ is purely combinatorial and independent of the coefficients:
\[
       (ds)_J=\sum_{a=0}^{|J|-1}(-1)^a s_{J\setminus\{j_a\}},
\tag{12.1}
\]
where a missing component is interpreted as zero.

\subsection{The three sign cases}

\begin{theorem}[Monomial summands]\label{thm:monomial-summands}
For a fixed exponent vector $e$:
\begin{enumerate}[label=\arabic*.]
\item If every $e_i\geq0$, then $C^\bullet(e)$ has cohomology $R$ in degree $0$ and zero in positive
      degrees.
\item If every $e_i<0$, then $C^\bullet(e)$ has a single nonzero term, in degree $n$, and its
      cohomology is $R$ there.
\item If $\operatorname{Neg}(e)$ is nonempty and proper, then $C^\bullet(e)$ is contractible by an
      explicit primitive-recursive homotopy.
\end{enumerate}
\end{theorem}

\begin{proof}
In the first case, every subset component is present.  This is the ordered cochain complex of a full
$n$-simplex.  Its degree-zero constant tuples give one copy of $R$, and the usual insertion homotopy
contracts the positive degrees.

In the second case, a subset containing all $n+1$ indices must be the full set, so the only term is in
degree $n$.

For the mixed case, choose the least index
$i_*\notin\operatorname{Neg}(e)$.  For a $p$-cochain $s$, define the $(p-1)$-cochain $hs$ on an
ordered set $I=\{i_0<\cdots<i_{p-1}\}$ by zero when $i_*\in I$, and otherwise by inserting $i_*$:
\begin{equation}\label{eq:insertion-homotopy}
 (hs)_I=(-1)^a s_{I\cup\{i_*\}},
 \qquad
 a=\#\set{i\in I:i<i_*}.
\end{equation}
The component on the right exists because adding $i_*$ does not change whether the set contains
$\operatorname{Neg}(e)$.  A finite sign computation pairs every deletion not involving $i_*$ once in
$dh$ and once with the opposite sign in $hd$; the remaining deletion of $i_*$ contributes $s_I$.
Hence
\[
                              dh+hd=\mathrm{id}.
\]
All sums have at most $n+1$ terms.
\end{proof}

\subsection{Cohomology table}

\begin{theorem}[Finite \v{C}ech cohomology of projective space]\label{thm:pn-cohomology}
For $n\geq1$, the standard-cover \v{C}ech cohomology is
\[
 \check H^q(\mathbf P^n_R,\cO(d))\cong
 \begin{cases}
   R[X_0,\ldots,X_n]_d,& q=0,\ d\geq0,\\[0.2em]
   \displaystyle
   \left(\frac1{X_0\cdots X_n}
   R[X_0^{-1},\ldots,X_n^{-1}]\right)_d,& q=n,\ d\leq-n-1,\\[0.5em]
   0,&\text{otherwise.}
 \end{cases}
\]
The two nonzero modules are finite free.  Their ranks are
\[
    \operatorname{rank}H^0=\binom{n+d}{n}\quad(d\geq0),
    \qquad
    \operatorname{rank}H^n=\binom{-d-1}{n}\quad(d\leq-n-1).
\]
The isomorphisms and zero contractions are primitive-recursive on finite cochain codes.
\end{theorem}

\begin{proof}
Decompose every finite cochain coefficientwise by exponent vector.  By
Theorem~\ref{thm:monomial-summands}, only all-nonnegative exponent vectors survive in degree zero and
all-negative vectors survive in degree $n$.  The all-nonnegative vectors of total degree $d$ are the
ordinary degree-$d$ monomials, which exist only for $d\geq0$ and are counted by
$\binom{n+d}{n}$.

For an all-negative vector write $e_i=-a_i-1$ with $a_i\geq0$.  Then
\[
                   \sum_i a_i=-d-(n+1).
\]
Such vectors exist exactly when $d\leq-n-1$, and stars-and-bars counts them by
$\binom{-d-1}{n}$.  The normal-form map keeps precisely those monomial coefficients and applies the
insertion homotopy to every mixed-sign summand.
\end{proof}

For $n=1$, Theorem~\ref{thm:pn-cohomology} specializes to
Theorems~\ref{thm:p1-h0} and~\ref{thm:p1-h1}.

\subsection{Primitive-recursive realization}

The implementation needs the following fixed operations:

\begin{itemize}
\item signed exponent-vector coding and bounded coordinate access;
\item the set $\operatorname{Neg}(e)$ as a bit mask;
\item enumeration of subsets of $\{0,\ldots,n\}$ by bit masks;
\item alternating-sign finite sums;
\item insertion of the least index outside a bit mask;
\item finite maps from exponent vectors to coefficient tuples;
\item coefficientwise splitting into all-nonnegative, all-negative, and mixed classes.
\end{itemize}

All are primitive recursive.  To prove \eqref{eq:insertion-homotopy}, an object-iPRA development can
avoid abstract matrices.  For each output component and each deletion index, a bounded case split
checks whether the deleted index is $i_*$, below it, or above it.  A finite accumulator records the
paired terms and their signs.  The induction formula is an equality of natural-number-coded finite
sums.

\subsection{Executable matrix check}

The companion program independently constructs the finite matrices of $d$ and $h$ for every proper
nonempty negative-index set in dimensions $1\leq n\leq6$.  It verifies
$dh+hd=I$ over the integers in 240 mixed-sign cases.  This is not a proof for all $n$, but it checks
the exact sign convention and the finite-index implementation used by the Lean plan.

\subsection{Conventional cohomology}

The standard cover has affine finite intersections, and $\cO(d)$ is quasi-coherent.  The conventional
comparison theorem therefore identifies the finite \v{C}ech calculation with
$H^q(\mathbf P^n_R,\cO(d))$.  The Stacks Project gives the same exponent-vector decomposition and
insertion homotopy.  The new foundational point is that the entire calculation itself is finite,
explicit, and admits a direct route to quantifier-free induction.

\section{Resolution map for the repository difficulties}

\subsection{What the new framework resolves}

The contribution is architectural rather than a claim that every difficult theorem in the
repository is now proved.  The following table separates resolved representation problems from
remaining mathematical producer problems.

\begin{longtable}{@{}p{0.27\textwidth}p{0.31\textwidth}p{0.34\textwidth}@{}}
\toprule
Difficulty & Resolution supplied here & Remaining boundary \\
\midrule
\endfirsthead
\toprule
Difficulty & Resolution supplied here & Remaining boundary \\
\midrule
\endhead
Quotients without decidable normal forms & Proof-relevant ideal-combination equality; operations transform witnesses. & Finding a witness may require a separate Gröbner or elimination theorem. \\
Localization with zero divisors & Equality carries a killing exponent; transitivity and congruence have explicit formulas. & Invertibility or annihilation witnesses must be supplied or produced. \\
Finite local-to-global gluing & Power-cover constructor plus explicit module gluing and uniqueness. & Arbitrary infinite covers are not internalized; use finite refinements or finite observations. \\
Scheme representation & Certified finite affine atlas with localization isomorphisms and cocycles. & No single completed internal universe of all schemes or prime ideals is claimed. \\
Sheaf condition & Exact equalizer for modules on finite principal covers. & General sheaf categories and derived functors require a larger coded library. \\
Projective geometry & Standard $\Proj$ charts and transitions are primitive-recursively generated. & Global projective existence theorems may still require nontrivial algebraic certificates. \\
Projective cohomology & Explicit finite \v{C}ech normal forms and contracting homotopy for $\cO(d)$. & Formal comparison with derived-functor cohomology is a separate extension. \\
Chart/branch nonuniqueness & Combine certified atlases with the prior finite-use relational framework; require invariant outputs. & Discontinuous exact selectors remain unavailable without stronger input. \\
Object-theory boundary & Exact list of \code{PR}/\code{PTm}/\code{Prf PRAAx} targets and quantifier-free invariants. & The final Lean implementation and compilation have not been completed in this archive. \\
\bottomrule
\end{longtable}

\subsection{Integration with Diophantine geometry}

The existing Diophantine packet represents rational projective points, heights, polynomial maps,
conics, local congruence obstructions, and finite descent.  The atlas layer extends that work in four
ways.

\begin{enumerate}[label=\arabic*.]
\item \textbf{Local equations.}  A projective variety can be dehomogenized on finitely many charts,
      and local algebraic manipulations can be glued or compared by localization witnesses.
\item \textbf{Rational maps with domains.}  A denominator-nonvanishing condition is represented by a
      principal open.  Several formulas for the same rational map are compared on overlaps rather
      than forced into one global syntactic expression.
\item \textbf{Line bundles and sections.}  Height constructions often use homogeneous forms or
      divisors.  Transition factors and local sections can now be represented and checked on a
      finite atlas.
\item \textbf{Local-to-global obstructions.}  Congruence, real, or $p$-adic local certificates can be
      attached to chart data, while exact algebraic descent witnesses handle the finite overlaps.
\end{enumerate}

A bounded rational-point search on a projective scheme proceeds as before: enumerate primitive
coordinate tuples up to a supplied height, check homogeneous equations, and quotient by projective
scaling.  The atlas does not create the height bound.  It provides the geometric operations needed
to state and verify local computations around the enumerated points.

\subsection{Faltings and Siegel}

The Faltings and Siegel packets correctly distinguish a finite cardinality statement from an
effective height cutoff.  The present framework sharpens the interface but does not remove that
obstruction.

A positive theorem can have the form:

\begin{quote}
Given a presented projective curve $X$, a height bound $B$, and a certificate proving that every
rational point has height at most $B$, the bounded enumeration returns a complete list; every local
chart computation and overlap equality is certified by the atlas.
\end{quote}

The bounded enumeration, chart conversion, and completeness-from-bound implication are primitive
recursive.  The production of $B$ for general high-genus curves is the missing arithmetic theorem.
A count bound, a finite-dimensional span bound, or a point-count estimate over finite fields is not a
substitute.

For Runge-type integral-point problems, explicit pole data and a supplied separation/height estimate
can be attached to the finite projective compactification.  This may simplify the exact bookkeeping
of points at infinity, divisors, and local functions, but the quantitative Runge inequality remains
the substantive producer.

\subsection{Mordell--Weil descent}

The Mordell--Weil packet has an explicit finite-height descent conditional on actual representatives
of the quotient by doubling.  A finite group scheme, isogeny, or Selmer computation can now be
presented chartwise with exact localization certificates.  This is useful for verifying the descent
maps and local images.

It still does not follow that a finite abstract Selmer class has a rational point representative.
The atlas framework can verify a supplied representative and its compatibility; it cannot create a
rational point from a purely cohomological class without an additional theorem.

\subsection{Puiseux expansions and singularities}

For exact polynomial coefficients, Newton polygons, finite truncations, and branch equations can be
encoded as finite algebraic certificates.  A singularity may be covered by finitely many candidate
charts or blow-up charts; relational finite-use composition can avoid an arbitrary global branch
label.

For arbitrary Cauchy coefficient names, the exact vanishing tests that distinguish singular strata
remain discontinuous.  The correct interfaces are:

\begin{itemize}
\item a supplied exact algebraic stratum or Newton-polygon certificate;
\item a derivative/separation witness giving a regular local branch;
\item finite root clusters and branch-invariant observations;
\item approximate solution germs with residual bounds.
\end{itemize}

The finite atlas does not turn approximation into exact equality.

\subsection{Analytic number theory packets}

Chebotarev, Siegel--Walfisz, Bombieri--Vinogradov, and related packets require uniform analytic
estimates, exceptional-zero control, field arithmetic, and explicit error terms.  Finite algebraic
geometry can organize number fields, finite covers, discriminants, and Frobenius data, but it does
not supply the missing estimates.  The no-universal-evaluator and finite-observation disciplines
continue to apply.

\subsection{Further algebraic-geometry targets}

The framework gives a clear route to several additional iPRA projects.

\begin{enumerate}[label=\arabic*.]
\item \textbf{Blow-ups of finitely presented affine schemes.}  Use the standard finite charts of the
      Rees algebra and localization transition maps.
\item \textbf{Grassmannians.}  Use Pl\"ucker coordinates, finitely many standard minors, and explicit
      transition matrices.
\item \textbf{Finite morphisms and normalization certificates.}  Present the target algebra as a
      finite module with multiplication tables and verify integral equations.
\item \textbf{B\'ezout in zero-dimensional certified intersections.}  Combine resultants or a
      homogeneous Gröbner certificate with finite quotient-rank calculations.
\item \textbf{Riemann--Roch for explicitly presented curves.}  First formalize divisor/line-bundle
      \v{C}ech complexes and finite-dimensional normal forms; any general theorem still needs
      quantitative algebraic inputs.
\item \textbf{Cohomology of toric varieties.}  Fans are finite combinatorial objects, affine toric
      charts are monoid algebras, and line-bundle cohomology often reduces to finite lattice-region
      calculations.
\item \textbf{Finite descent for vector bundles.}  Transition matrices and principal-cover gluing
      already provide the base layer for explicit projective-module calculations.
\end{enumerate}

The common theme is that the geometric theorem should expose a finite atlas, a finite complex, or a
finite certificate rather than conceal infinitary semantics behind an unexamined object code.

\section{Lean formalization: exact definitions and a route to \texorpdfstring{$\Prf\,\PRAAx$}{Prf PRAAx}}

\subsection{The formal endpoint}

The repository distinguishes three formal systems that must not be conflated.  Lean's ambient type
 theory is used to define syntax and metatheorems; the ordinary first-order presentation of iPRA has
primitive-recursive terms \code{PTm}, axioms \code{PRAAx}, and derivations \code{Prf}; and the
smaller native scalar core has its own derivation relation.  The checked bridge is of the form
\begin{lstlisting}[language={}] 
theorem pra_iff_native
    {Gamma : List (Fm PTm n)} {phi : Fm PTm n} :
  (Gamma |- [PRAAx] phi) <->
    NPrf (Gamma.map NFm.ofPRA) (NFm.ofPRA phi)
\end{lstlisting}
Consequently a host theorem such as
\begin{lstlisting}[language={}] 
theorem power_cover_semantic : ...
\end{lstlisting}
is not yet an iPRA formalization.  The relevant endpoint is a term such as
\begin{lstlisting}[language={}] 
def powerCoverObjectProof :
    Prf PRAAx [] powerCoverFormula := ...
\end{lstlisting}
and only afterwards may \code{pra_iff_native} transport it to the native calculus.

The source files in this archive are blueprints.  No Lean executable or complete repository
dependency closure was available in the preparation environment, so they are not claimed to
compile.  They contain no \code{axiom}, \code{sorry}, or \code{admit}; unfinished results are named
\emph{targets}, not asserted theorems.

\subsection{Two formalization layers}

The proposed development has two deliberately separate layers.

\paragraph{Host certificate layer.}
A \code{CertSetoid} contains numerical codes, a Boolean validity predicate, a proof-relevant
relation \code{eqW x y w}, and explicit witness transformers for reflexivity, symmetry, and
transitivity.  A certified map carries equality witnesses to equality witnesses.  Certified rings
and modules add congruence transformers for their operations.  This layer is useful for testing the
interfaces and for proving generic metatheorems, but its quantifiers live in Lean rather than in the
object theory.

\paragraph{Object-iPRA layer.}
Every numerical operation is replaced by one fixed \code{PR} code.  Its application is a
\code{PTm}; every checker and invariant is an \code{Fm PTm n}; and all correctness theorems are
values of \code{Prf PRAAx}.  A theorem schema may fix the definitions of a presented ring or module
externally and remain uniform in ordinary numerical parameters.  It must not evaluate arbitrary
codes for primitive-recursive programs, since there is no primitive-recursive universal evaluator
for all such codes.

\subsection{Codecs and totalization}

All objects of the internal theory are natural numbers.  The first module should therefore supply
fixed primitive-recursive codecs for:
\begin{itemize}
\item signed integers and finite tuples;
\item lists with length, bounded lookup, update, append, map, fold, maximum, and least bounded search;
\item exponent vectors, monomials, sparse polynomials, and Laurent polynomials;
\item finite matrices and bit-mask subsets of $\{0,\ldots,n\}$;
\item ideal-combination witnesses, localization witnesses, cover certificates, and descent data.
\end{itemize}
Every operation is total.  Malformed input receives a default output together with a failure flag.
Theorems assume the finite validity predicate, so no partial function is hidden in the object
language.  The basic reconstruction lemmas have the shape
\[
  \Valid(x)=1\;\longrightarrow\;
  \operatorname{decode}(\operatorname{encode}(x))=x,
\]
and bounded access lemmas are proved by quantifier-free induction on the list prefix.

\subsection{Certified-setoid library}

For each presentation $A$, introduce fixed primitive-recursive predicates
\[
        \Valid_A(x),\qquad \EqW_A(x,y,w).
\]
The object library should construct the following derivations, uniformly in the displayed codes:
\begin{align*}
\Valid_A(x)&\to \EqW_A(x,x,\mathsf{reflW}(x)),\\
\EqW_A(x,y,w)&\to \EqW_A(y,x,\mathsf{symmW}(x,y,w)),\\
\EqW_A(x,y,u)\wedge\EqW_A(y,z,v)
 &\to \EqW_A(x,z,\mathsf{transW}(x,y,z,u,v)).
\end{align*}
For every operation $F$, add a witness transformer
\[
   \EqW_A(x,x',u)\wedge\EqW_B(y,y',v)
   \to \EqW_C(F(x,y),F(x',y'),\mathsf{mapW}(u,v)).
\]
The library never asserts $\exists w\,\EqW_A(x,y,w)$ merely because two semantic quotient classes
are equal in Lean.  The witness is part of the finite mathematical data.

A useful generic lemma is bounded-fold congruence.  Suppose two lists have the same length and each
pair of corresponding entries comes with an equality witness.  Quantifier-free induction on the
prefix produces a witness between the two additive or multiplicative folds.  This lemma will be
used repeatedly in ideal identities, multinomial expansions, and \v{C}ech differentials.

\subsection{Presented quotient rings}

For
\[
 A=\Z[X_1,\ldots,X_s]/(g_0,\ldots,g_{r-1}),
\]
a code for an element is a polynomial code.  A witness that $p=q$ is a tuple
$(h_0,\ldots,h_{r-1})$ satisfying the coefficientwise identity
\[
                     p-q=\sum_{i<r}h_i g_i.
\]
The checker expands every product and compares a bounded coefficient range determined from the
input supports.  The soundness proof uses only finite polynomial arithmetic.

The witness transformers are explicit.  Negation negates every multiplier; addition concatenates
or adds multiplier lists; and multiplication uses
\[
 pq-p'q'=(p-p')q+p'(q-q').
\]
Associativity, commutativity, distributivity, and the unit laws can be witnessed by the zero
multiplier list once the corresponding polynomial identity is normalized coefficientwise.  No
Gr\"obner basis is needed for soundness.  A Gr\"obner or elimination algorithm is a separate
\emph{producer} that may find witnesses in useful subclasses.

The principal object theorem should be exposed as a derivation of a formula asserting:
\begin{quote}
if the ideal-certificate checker accepts, then the represented quotient elements are equal, and
all quotient operations preserve accepted equality certificates.
\end{quote}

\subsection{Localization witnesses}

Fix a presented ring $A$ and an element $f$.  A fraction is a pair $(a,n)$ representing $a/f^n$.
Equality of $(a,n)$ and $(b,m)$ is witnessed by a number $k$ and a base-ring witness for
\[
                       f^k(f^m a-f^n b)=0.                 \tag{14.1}\label{eq:loc-witness-object}
\]
The object formalization should implement the following witness algebra.

\begin{itemize}
\item Reflexivity uses $k=0$ and the zero witness.
\item Symmetry negates the witness in \eqref{eq:loc-witness-object}.
\item For transitivity, suppose $k$ kills
      $f^m a-f^n b$ and $\ell$ kills $f^r b-f^m c$.  Multiplying the first
      relation by $f^{\ell+r}$ and the second by $f^{k+n}$ and adding cancels
      the middle term.  The result supplies a safe killing exponent
      $k+\ell+m$ for $f^r a-f^n c$; a slightly smaller exponent can be used
      after a more exact bookkeeping proof.
\item Addition and multiplication expand the cross differences and combine the supplied witnesses
      by finite ring identities.
\end{itemize}

Every exponentiation is a fixed primitive recursion.  The loop invariant records the current power,
the accumulated product, and an accepted equality witness.  The induction formula is
quantifier-free.

\subsection{The power-cover constructor}

The most important generic producer starts with a checked identity
\[
                         \sum_{i<r}c_i f_i=1               \tag{14.2}\label{eq:unit-cover-object}
\]
and a number $N\geq1$.  Set $M=r(N-1)+1$.  Enumerate every weak composition
$\alpha\in\N^r$ of $M$ and expand $(\sum c_i f_i)^M$.  The pigeonhole bound gives an index
$i$ with $\alpha_i\geq N$; choose the least such index and assign that monomial to the accumulator
$d_i$ after factoring $f_i^N$.  The output satisfies
\[
                         \sum_{i<r}d_i f_i^N=1.             \tag{14.3}\label{eq:power-cover-object}
\]

A direct object implementation uses nested bounded loops rather than an abstract multinomial
polynomial.  Its state contains:
\begin{enumerate}[label=(\roman*)]
\item the current exponent vector and its remaining sum;
\item the partial multinomial coefficient;
\item the chosen least large coordinate;
\item the coefficient accumulators $d_0,\ldots,d_{r-1}$;
\item the partial expansion value and the partial factored value.
\end{enumerate}
The invariant says that the two partial values are equal and that every processed composition was
assigned to exactly one accumulator.  All quantifiers are bounded folds represented by Boolean
prefix predicates.  A final finite-power calculation rewrites the expansion by \eqref{eq:unit-cover-object} to
$1^M=1$.

The endpoint should be
\begin{lstlisting}[language={}] 
def powerCoverObjectProof :
  Prf PRAAx [] powerCoverFormula := ...
\end{lstlisting}
where \code{powerCoverFormula} states validity of the input certificate, $N\geq1$, and acceptance
of the output identity checker.

\subsection{Principal-cover uniqueness and gluing}

For joint monicity, the input contains for each $i<r$ a local annihilation exponent $k_i$ and a
witness for $f_i^{k_i}m=0$.  Let $K=\max_i k_i$.  Raise each equation to the common exponent,
construct $d_i$ with $\sum d_i f_i^K=1$, and calculate
\[
                    m=\sum_i d_i f_i^K m=0.
\]
The induction invariant for the last fold records a witness that the accumulated partial sum equals
zero.

For existence, first homogenize the denominators of local sections to a common exponent $N$.
Take a maximum $K$ of the overlap-killing exponents and a safe $q=N+K+1$.  Form numerators $b_i$
so that the accepted overlap certificates imply
\[
                         f_j^q b_i=f_i^q b_j                \tag{14.4}
\]
for every $i,j$.  Construct $d_i$ with $\sum_i d_i f_i^q=1$ and output
\[
                              m=\sum_i d_i b_i.              \tag{14.5}\label{eq:glue-formula-object}
\]
For a fixed $j$, a bounded fold proves
\[
 f_j^q m=\sum_i d_i f_j^q b_i
          =\sum_i d_i f_i^q b_j=b_j.
\]
This equation is itself the localization equality witness between $m/1$ and the given local
section on $D(f_j)$.  A double-loop invariant covers all pairs $(i,j)$; no induction formula says
``for every localization''.

The two final statements should have types analogous to
\begin{lstlisting}[language={}] 
def principalCoverUniqueObjectProof :
  Prf PRAAx [] principalCoverUniqueFormula := ...

def principalCoverGlueObjectProof :
  Prf PRAAx [] principalCoverGlueFormula := ...
\end{lstlisting}

\subsection{Finite atlases and quasi-coherent data}

A finite affine atlas is a natural-number code containing:
\begin{itemize}
\item finitely presented chart rings $A_i$;
\item for each ordered pair $(i,j)$, finite principal-open refinements of the overlap;
\item forward and inverse localization homomorphisms with accepted composition witnesses;
\item triple-overlap cocycle witnesses;
\item a finite cover certificate wherever a family of principal opens is asserted to cover.
\end{itemize}
A quasi-coherent module contains chart modules $M_i$ and transition isomorphisms after the same
localizations.  Validity is one Boolean fold over the finite atlas.  The gluing theorem above gives
the equalizer law on every chart, and finite-family composition gives the corresponding law on the
atlas.  Morphisms, kernels, finite direct sums, and tensor products can then be added by certified
module operations.

This layer should be formalized as a library of certificate transformations, not as an internal
quantification over prime ideals or over all open subsets.  A separate semantic theorem in Lean may
interpret a valid atlas as a conventional glued scheme.

\subsection{The projective-line modules}

The $\mathbf P^1$ development uses sparse Laurent-polynomial codes.  Define
\[
       \delta_d(a,b)=a(t)-t^d b(t^{-1}).
\]
For a cocycle $c=\sum_e c_e t^e$, one bounded scan places:
\begin{align*}
 e\geq0 &\quad\hbox{in the first chart polynomial }a,\\
 e\leq d&\quad\hbox{in the second chart polynomial }b,\\
 d<e<0&\quad\hbox{in the obstruction }h.
\end{align*}
The invariant after processing a prefix of the finite support is
\[
       c_{<s}=\delta_d(a_s,b_s)+h_s,
       \qquad \operatorname{supp}(h_s)\subseteq\set{e:d<e<0}.
\]
This yields an object proof of \code{p1LaurentSplitFormula}.  Uniqueness is coefficientwise: neither
chart coboundary contains a monomial in the middle interval.  The $H^0$ theorem uses the same support
separation on a compatible pair.  In particular, at $d=0$ the global functions on $\mathbf P^1_R$
are exactly the constants in $R$.

\subsection{The projective-space homotopy}

Encode a subset $I\subseteq\{0,\ldots,n\}$ as a bit mask and an exponent vector as a signed list.
For one exponent vector $e$, let $N(e)$ be the negative-coordinate mask.  A component occurs at
$I$ exactly when $N(e)\subseteq I$.  The differential is a bounded alternating sum over deleted
bits.

When $N(e)$ is nonempty and not full, choose the least $i_*\notin N(e)$.  The homotopy inserts that
bit with sign determined by the number of smaller elements.  To prove
\[
                              dh+hd=\mathrm{id},             \tag{14.6}\label{eq:cech-homotopy-object}
\]
fix one output mask and scan all deletion positions.  The state records the unmatched term, the
paired contribution below $i_*$, and the paired contribution above $i_*$.  Terms not deleting
$i_*$ occur twice with opposite parity; the deletion of $i_*$ is the unique identity contribution.
The parity and bit-mask cases are decidable numerical predicates, so the induction is
quantifier-free.

The all-nonnegative case is the full simplex complex: the same insertion homotopy contracts the
positive degrees, while a coefficient projection identifies degree-zero constants.  The all-negative
case contains only the full mask.  A finite scan over the input cochain's support applies the
appropriate normalizer to each exponent vector.  This gives the object formula
\code{pnOdCechFormula} and the binomial-rank calculations by bounded counting of weak
compositions.

\subsection{Suggested module order}

A practical repository integration order is:
\begin{longtable}{@{}p{0.27\textwidth}p{0.65\textwidth}@{}}
\toprule
Module & Main contents and endpoint \\
\midrule
\endfirsthead
\toprule
Module & Main contents and endpoint \\
\midrule
\endhead
\code{AlgebraicGeometry/Codecs.lean} & Signed/list/polynomial/Laurent/matrix/subset codes and their PR graphs. \\
\code{AlgebraicGeometry/CertSetoid.lean} & \code{Valid}, \code{EqW}, witness transformers, bounded-fold congruence. \\
\code{AlgebraicGeometry/PresentedRing.lean} & Ideal-combination checker and quotient-operation witnesses. \\
\code{AlgebraicGeometry/Localization.lean} & Fraction codes and localization equivalence/congruence proofs. \\
\code{AlgebraicGeometry/PowerCover.lean} & Weak compositions, multinomial fold, and \code{powerCoverObjectProof}. \\
\code{AlgebraicGeometry/PrincipalDescent.lean} & Joint monicity and \code{principalCoverGlueObjectProof}. \\
\code{AlgebraicGeometry/Atlas.lean} & Finite affine atlases, refinements, morphisms, and module descent data. \\
\code{AlgebraicGeometry/ProjectiveLine.lean} & Two-chart atlas and Laurent normal forms for every $\cO(d)$. \\
\code{AlgebraicGeometry/ProjectiveCech.lean} & Exponent decomposition, insertion homotopy, and projective-space table. \\
\code{AlgebraicGeometry/NativeBridge.lean} & Applications of \code{pra_iff_native} to already constructed object derivations. \\
\bottomrule
\end{longtable}

\subsection{Completion criteria}

The formalization should be called complete only when all of the following hold:
\begin{enumerate}[label=\arabic*.]
\item the modules compile against the current \code{ConciseIPRA} sources;
\item the advertised formulas use the actual repository syntax and not a surrogate proof relation;
\item the principal theorems are inhabited \code{Prf PRAAx} values;
\item an axiom audit reports no \code{sorryAx}, custom axiom, or hidden native decision oracle;
\item finite executable tests agree with the object-level producers on a representative corpus;
\item the native translations are performed only after the preceding object proofs exist.
\end{enumerate}
The present archive supplies the mathematical derivations, numerical algorithms, host interfaces,
and this dependency order, but does not claim these completion criteria have already been met.

\section{Executable validation, theorem status, and reproducibility}

\subsection{Purpose of the executable companion}

The file \code{code/zariski\_descent.py} is an exact-arithmetic reference implementation of the
central finite algorithms.  It uses only the Python standard library and never uses floating point.
It is not an iPRA proof checker and does not replace the universal arguments in the preceding
sections.  Its role is to check formulas, sign conventions, exponent bookkeeping, and examples
independently of the LaTeX derivation.

The implementation includes:
\begin{itemize}
\item a minimal exact commutative-ring interface, with instances for $\Z$ and $\Z/m\Z$;
\item verification and exponentiation of unit-ideal cover certificates;
\item the multinomial power-cover coefficient constructor;
\item principal-cover gluing for the module $M=A$;
\item sparse univariate polynomial arithmetic and ideal-certificate checking;
\item the $\mathbf P^1$ Laurent \v{C}ech differential and normal form;
\item projective-space cohomology rank formulas;
\item integer matrices for the monomial-wise differential and insertion homotopy.
\end{itemize}

\subsection{Recorded validation}

Running
\begin{lstlisting}[language={}] 
python3 -m unittest discover -s code -v
python3 code/zariski_descent.py
\end{lstlisting}
produced the log in \code{verification/python\_tests.txt} and the structured records in
\code{metadata/validation.json} and \code{verification/demo\_output.json}.  The eight unit tests
passed.  The recorded checks comprise:

\begin{longtable}{@{}p{0.32\textwidth}p{0.58\textwidth}@{}}
\toprule
Component & Recorded check \\
\midrule
\endfirsthead
\toprule
Component & Recorded check \\
\midrule
\endhead
Power covers & Eleven exact cases over $\Z$ and $\Z/12\Z$, including exponents larger than one. \\
Integer gluing & A two-chart example with unequal denominators recovers the global section $17$; the optimized common exponent is $2$, and the computed power-cover coefficients are $(7,-3)$. \\
Zero-divisor gluing & A cover in $\Z/12\Z$ recovers the global section $5$, demonstrating that the algorithm does not cancel nonunits or assume a domain. \\
Ideal certificate & An exact polynomial identity of the form $p-q=hg$ is accepted coefficientwise; a perturbed certificate is rejected. \\
$\mathbf P^1$ normal form & $1{,}200$ randomly generated Laurent cocycles satisfy exact reconstruction, and every obstruction term lies in $d<e<0$. \\
$\mathbf P^1$ global sections & Generated coefficient vectors give compatible chart pairs for several positive degrees and the expected empty basis for negative degrees. \\
Projective ranks & The code checks the closed rank formulas against direct bounded enumeration in the tested range. \\
$\mathbf P^n$ homotopy & For all proper nonempty negative-coordinate masks in dimensions $1\leq n\leq6$, $240$ matrix cases satisfy $dh+hd=I$ over $\Z$. \\
\bottomrule
\end{longtable}

The report uses the uniformly safe gluing exponent $q=N+K+1$.  The executable companion uses the
slightly sharper $q=\max(1,N+K)$ and then verifies the resulting cross identities.  This difference
is intentional: the paper favors a simple bound with slack, while the program demonstrates that the
slack is not mathematically essential.

\subsection{Worked principal-cover examples}

\paragraph{An integer example.}
Take the cover elements $f_0=2$, $f_1=3$ with certificate
\[
                         (-1)\cdot2+1\cdot3=1.
\]
Choose local fractions obtained from the global integer $17$ but with different denominator
exponents.  The compatibility checker clears the denominators, the gluing routine constructs a
common power cover, and the output is exactly $17$.  This test exercises denominator homogenization
rather than merely gluing two identical numerators.

\paragraph{A zero-divisor example.}
In $A=\Z/12\Z$, the elements $3$ and $4$ satisfy
\[
                         3\cdot3+1\cdot4=1\pmod{12}.
\]
Neither chart element is a unit, and $A$ has zero divisors.  Localizations and overlap equalities are
therefore genuinely proof-relevant.  The supplied local data glue to $5\in A$.  Every final
restriction equation is checked in the ring.  This example is a useful guard against accidentally
using cancellation in the human or Lean proof.

\subsection{Claim ledger}

The machine-readable file \code{metadata/claims.json} and the human-readable
\code{metadata/THEOREM\_LEDGER.md} distinguish the four relevant levels of completion.  The main
status is summarized here.

\begin{longtable}{@{}p{0.12\textwidth}p{0.48\textwidth}p{0.31\textwidth}@{}}
\toprule
ID & Result & Status in this archive \\
\midrule
\endfirsthead
\toprule
ID & Result & Status in this archive \\
\midrule
\endhead
F1 & Certified-setoid and congruence discipline & Full human proof and Lean interface; object derivation pending. \\
A1 & Quotient equality by ideal combinations & Full human proof and exact checker; object derivation pending. \\
A2 & Localization equivalence and ring congruence & Explicit witness formulas proved; Lean target supplied. \\
Z1 & Power-cover constructor & Full proof, exact implementation, finite tests; \code{Prf PRAAx} pending. \\
Z2 & Joint monicity on a principal cover & Full proof; object derivation pending. \\
Z3 & Principal-cover module gluing & Full proof including zero divisors; exact implementation for $M=A$; object derivation pending. \\
G1--G2 & Finite affine atlases and quasi-coherent descent data & Precise certificate semantics and constructions; no complete category-of-schemes library claimed. \\
P1--P3 & $\mathbf P^1$, $\cO(d)$, $H^0$, and $H^1$ & Full finite proofs and exact normal-form tests; object derivations pending. \\
PN & Mixed-sign contraction and the standard $\mathbf P^n$ table & Full combinatorial proof and finite matrix checks; object derivation pending. \\
C1 & Agreement with conventional sheaf cohomology & Uses the standard affine-cover comparison theorem; no internal derived-functor formalization claimed. \\
L1 & Lean completion route & Source-level blueprints and dependency order; not compiler-checked. \\
\bottomrule
\end{longtable}

\subsection{Negative checks and nonclaims}

The validation deliberately does not infer an unbounded theorem from a successful finite test.
Specifically:
\begin{itemize}
\item testing finitely many rings does not prove the power-cover theorem for every presented ring;
\item a bounded point search does not prove absence of points above its cutoff;
\item a successful certificate checker does not guarantee that a certificate producer exists for
      every true semantic statement;
\item native Lean propositions, once added, would not by themselves be \code{Prf PRAAx} values;
\item the standard \v{C}ech computation is not by itself a formalization of all derived-functor
      sheaf cohomology inside iPRA.
\end{itemize}

No claim is made here of general effective Siegel or Faltings, uniform Mordell--Weil generators,
exact singular Puiseux selection from arbitrary Cauchy coefficients, decidable equality of arbitrary
represented reals, or a primitive-recursive universal evaluator.

\subsection{Rebuilding the package}

From the extracted package root, the intended commands are:
\begin{lstlisting}[language={}] 
python3 -m unittest discover -s code -v
python3 code/zariski_descent.py > metadata/validation.json
cd paper
latexmk -pdf -interaction=nonstopmode -halt-on-error \
  ipra_algebraic_geometry_framework.tex
\end{lstlisting}
The included \code{build.sh} performs the tests, rebuilds the PDF, and writes an integrity manifest.
A normal TeX Live installation with the packages named in the master source is sufficient.

The upstream repository audit is pinned to
\code{bce5113f44d613f5625c274128ee761f35486fc5}.  At the final check, \code{main} was
\code{a2f72b0cdf57f8a3fefd34b272c3d471385e819c}; the Git object identifiers of \code{ConciseIPRA}, the iPRA
\code{README.md}, and \code{docs} were identical at the two commits.  The archive neither
modifies nor rebuilds the repository.  Its source audit is in
\code{metadata/SOURCE\_AUDIT.md}; its explicit limitations are in
\code{metadata/LIMITATIONS.md}.

\subsection{Archive inventory}

\begin{longtable}{@{}p{0.35\textwidth}p{0.55\textwidth}@{}}
\toprule
Path & Purpose \\
\midrule
\endfirsthead
\toprule
Path & Purpose \\
\midrule
\endhead
\code{paper/ipra\_algebraic\_geometry\_framework.tex} & Master editable LaTeX source. \\
\code{paper/sections/} & The complete section sources, including proofs and formalization plan. \\
\code{paper/ipra\_algebraic\_geometry\_framework.pdf} & Compiled report. \\
\code{code/zariski\_descent.py} & Exact finite reference algorithms. \\
\code{code/test\_zariski\_descent.py} & Unit and randomized regression tests. \\
\code{lean/*.lean} & Host structures, target propositions, and object-theory endpoint blueprint. \\
\code{lean/FORMALIZATION\_PLAN.md} & Dependency-ordered path to complete Lean and object-iPRA proofs. \\
\code{metadata/SOURCE\_AUDIT.md} & Pinned repository and external-source audit. \\
\code{metadata/THEOREM\_LEDGER.md} & Result-by-result proof status. \\
\code{metadata/claims.json} & Machine-readable claim levels and nonclaims. \\
\code{metadata/validation.json} & Exact executable validation record. \\
\code{verification/} & Test logs, demo output, PDF build report, and visual QA notes. \\
\code{MANIFEST.sha256} & Integrity hashes of delivered files other than the manifest itself. \\
\bottomrule
\end{longtable}

\section{References and source provenance}

The mathematical proofs in this report are written out rather than delegated to the references.
The references locate the conventional statements, the inspected repository interfaces, and related
constructive or formal developments.

\begin{thebibliography}{99}

\bibitem{proveit-pra}
V. Reshetnikov,
\emph{ProveIt: Primitive recursive arithmetic},
repository directory \url{https://github.com/VladimirReshetnikov/ProveIt/tree/bce5113f44d613f5625c274128ee761f35486fc5/Logic/PrimitiveRecursiveArithmetic},
inspected 15 November 2026 (Pacific time).

\bibitem{concise-ipra}
V. Reshetnikov,
\emph{Concise presentations of intuitionistic primitive recursive arithmetic},
\url{https://github.com/VladimirReshetnikov/ProveIt/tree/bce5113f44d613f5625c274128ee761f35486fc5/Logic/PrimitiveRecursiveArithmetic/ConciseIPRA}.
In particular, the Lean modules \code{PR.lean} and \code{NativeBiInterpretation.lean} define
\code{PR}, \code{PTm}, \code{PRAAx}, \code{Prf}, and \code{pra\_iff\_native}.

\bibitem{finite-use}
\emph{Demand-Driven Relational Certificates for Mathematics in iPRA},
ProveIt packet \code{ipra\_finite\_use\_framework\_2026-11-15},
\url{https://github.com/VladimirReshetnikov/ProveIt/tree/bce5113f44d613f5625c274128ee761f35486fc5/Logic/PrimitiveRecursiveArithmetic/docs/proofs/foundations/ipra_finite_use_framework_2026-11-15}.

\bibitem{coherent-methods}
\emph{Coherent finite methods for intuitionistic primitive recursive arithmetic},
ProveIt packet \code{ipra\_coherent\_methods\_2026-11-15},
\url{https://github.com/VladimirReshetnikov/ProveIt/tree/bce5113f44d613f5625c274128ee761f35486fc5/Logic/PrimitiveRecursiveArithmetic/docs/proofs/foundations/ipra_coherent_methods_2026-11-15}.

\bibitem{embedding-framework}
\emph{A general embedding framework for mathematics in iPRA},
ProveIt packet \code{ipra\_embedding\_framework},
\url{https://github.com/VladimirReshetnikov/ProveIt/tree/bce5113f44d613f5625c274128ee761f35486fc5/Logic/PrimitiveRecursiveArithmetic/docs/proofs/foundations/ipra_embedding_framework}.

\bibitem{local-global}
\emph{Local-to-global constructions in intuitionistic primitive recursive arithmetic},
ProveIt packet \code{ipra\_local\_global},
\url{https://github.com/VladimirReshetnikov/ProveIt/tree/bce5113f44d613f5625c274128ee761f35486fc5/Logic/PrimitiveRecursiveArithmetic/docs/proofs/foundations/ipra_local_global}.

\bibitem{stable-observables}
\emph{Stable observables from finite witnesses},
ProveIt packet \code{ipra\_stable\_observables},
\url{https://github.com/VladimirReshetnikov/ProveIt/tree/bce5113f44d613f5625c274128ee761f35486fc5/Logic/PrimitiveRecursiveArithmetic/docs/proofs/algebra/ipra_stable_observables}.

\bibitem{diophantine-geometry}
\emph{Diophantine geometry in intuitionistic primitive recursive arithmetic},
ProveIt packet \code{diophantine\_geometry\_ipra},
\url{https://github.com/VladimirReshetnikov/ProveIt/tree/bce5113f44d613f5625c274128ee761f35486fc5/Logic/PrimitiveRecursiveArithmetic/docs/proofs/diophantine-number-theory/diophantine_geometry_ipra}.

\bibitem{faltings-packet}
\emph{Faltings' theorem and intuitionistic PRA},
ProveIt packet \code{faltings\_ipra},
\url{https://github.com/VladimirReshetnikov/ProveIt/tree/bce5113f44d613f5625c274128ee761f35486fc5/Logic/PrimitiveRecursiveArithmetic/docs/proofs/diophantine-number-theory/faltings_ipra}.

\bibitem{siegel-packet}
\emph{Siegel's theorem and intuitionistic primitive recursive arithmetic},
ProveIt packet \code{siegel\_ipra},
\url{https://github.com/VladimirReshetnikov/ProveIt/tree/bce5113f44d613f5625c274128ee761f35486fc5/Logic/PrimitiveRecursiveArithmetic/docs/proofs/diophantine-number-theory/siegel_ipra}.

\bibitem{mordell-weil-packet}
\emph{Constructive Mordell--Weil descent in iPRA},
ProveIt packet \code{mordell\_weil\_ipra},
\url{https://github.com/VladimirReshetnikov/ProveIt/tree/bce5113f44d613f5625c274128ee761f35486fc5/Logic/PrimitiveRecursiveArithmetic/docs/proofs/diophantine-number-theory/mordell_weil_ipra}.

\bibitem{stacks-localization}
The Stacks Project Authors,
\emph{Localization}, Tag \code{00CM},
\url{https://stacks.math.columbia.edu/tag/00CM}.

\bibitem{stacks-affine}
The Stacks Project Authors,
\emph{Affine schemes}, Tag \code{01HR},
\url{https://stacks.math.columbia.edu/tag/01HR}.

\bibitem{stacks-principal-open}
The Stacks Project Authors,
\emph{Affine schemes and standard opens}, Tag \code{01HR}, and
\emph{The category of affine schemes}, Tag \code{01HX},
\url{https://stacks.math.columbia.edu/tag/01HR} and
\url{https://stacks.math.columbia.edu/tag/01HX}.

\bibitem{stacks-qc}
The Stacks Project Authors,
\emph{Quasi-coherent sheaves on affines}, Tag \code{01I6},
\url{https://stacks.math.columbia.edu/tag/01I6}.

\bibitem{stacks-gluing}
The Stacks Project Authors,
\emph{Glueing schemes}, Tag \code{01JA},
\url{https://stacks.math.columbia.edu/tag/01JA}.

\bibitem{stacks-cech}
The Stacks Project Authors,
\emph{\v{C}ech cohomology and cohomology}, Tag \code{01EO},
\url{https://stacks.math.columbia.edu/tag/01EO}.

\bibitem{stacks-projective-cohomology}
The Stacks Project Authors,
\emph{Cohomology of projective space}, Tag \code{01XS},
\url{https://stacks.math.columbia.edu/tag/01XS}.

\bibitem{buchberger}
B. Buchberger,
\emph{An Algorithm for Finding the Basis Elements of the Residue Class Ring of a
Zero Dimensional Polynomial Ideal},
Ph.D. thesis, University of Innsbruck, 1965; English translation, 2006.

\bibitem{coquand-lombardi}
T. Coquand and H. Lombardi,
\emph{A logical approach to abstract algebra},
Mathematical Structures in Computer Science 16 (2006), 885--900.

\bibitem{zhang-proj}
J. Zhang,
\emph{Formalising the Proj Construction in Lean},
Interactive Theorem Proving (ITP 2023), LIPIcs 268, Article 35.

\bibitem{hartshorne}
R. Hartshorne,
\emph{Algebraic Geometry},
Graduate Texts in Mathematics 52, Springer, 1977.

\bibitem{wiki-ag}
Wikipedia contributors,
\emph{Algebraic geometry},
\url{https://en.wikipedia.org/wiki/Algebraic_geometry},
used only as the user-specified orientation checklist.

\end{thebibliography}

\subsection*{Provenance note}

A more detailed path-by-path audit, including what was and was not reused, is included as
\code{metadata/SOURCE\_AUDIT.md}.  The external references describe conventional algebraic
geometry and related constructive methods; they are not substitutes for the finite arithmetic
proofs presented in this report.


\end{document}
